% Emacs, this is -*-latex-*-

\ex{Convers\~ao para grafos de fatores}
\label{ex:DGM-vs-UGM-alt}
\begin{exenumerate}
    
    \item Desenhe um grafo n\~ao direcionado e um grafo de fatores n\~ao direcionado para
    $p(x_1,x_2,x_3) = p(x_1) p(x_2) p(x_3 | x_1,x_2)$
    
    \begin{solution}
        \begin{center}
            \scalebox{0.9}{ % x,y
                \begin{tikzpicture}[ugraph]
                    \node[cont] (x1) at (0,2) {$x_1$};
                    \node[cont] (x2) at (2,2) {$x_2$};
                    \node[cont] (x3) at (1,-0.2) {$x_3$};
                    
                    \draw(x1) -- (x2);
                    \draw(x2) -- (x3);
                    \draw(x3) -- (x1);
                \end{tikzpicture}
                \begin{tikzpicture}[ugraph]
                    \node[cont] (x1) at (0,2) {$x_1$};
                    \node[cont] (x2) at (2,2) {$x_2$};
                    \node[fact, label=right: $p(x_3 | x_1\, x_2)$] (f3) at (1,1) {};
                    \node[cont] (x3) at (1,-0.2) {$x_3$};
                    \node[fact, label=left: $p(x_1)$] (f1) at (0,3) {};
                    \node[fact, label=right: $p(x_2)$] (f2) at (2,3) {};
                    
                    \draw(x1) -- (f3);
                    \draw(x2) -- (f3);
                    \draw(f3) -- (x3);
                    \draw[-](f1) -- (x1);
                    \draw[-](f2) -- (x2);
                \end{tikzpicture}
            }
        \end{center}
    \end{solution}
    
    \item Desenhe um grafo de fatores n\~ao direcionado para o modelo gr\'afico direcionado definido pelo grafo abaixo.
    
    \begin{center}
        \scalebox{0.9}{
            \begin{tikzpicture}[dgraph]
                \node[cont] (y1) at (0,0) {$y_1$};
                \node[cont] (y2) at (2,0) {$y_2$};
                \node[cont] (y3) at (4,0) {$y_3$};
                \node[cont] (y4) at (6,0) {$y_4$};
                \node[cont] (x1) at (0,2) {$x_1$};
                \node[cont] (x2) at (2,2) {$x_2$};
                \node[cont] (x3) at (4,2) {$x_3$};
                \node[cont] (x4) at (6,2) {$x_4$};
                \draw(x1)--(y1);\draw(x2)--(y2);\draw(x3)--(y3);\draw(x4)--(y4);
                \draw(x1)--(x2);\draw(x2)--(x3);\draw(x3)--(x4);
        \end{tikzpicture}}
    \end{center}
    
    \begin{solution}
        O grafo especifica modelos probabil\'isticos que se fatorizam como
        $$p(x_1, \ldots, x_4, y_1, \ldots, y_4) = p(x_1) p(y_1 | x_1)
        \prod_{i=2}^4 p(y_i |x_i) p(x_i | x_{i-1})$$
        \'E o grafo para um modelo oculto de Markov (HMM). O grafo de fatores correspondente \'e mostrado abaixo.
        \begin{center}
            \scalebox{0.9}{
                \begin{tikzpicture}[ugraph]
                    \node[cont] (y1) at (0,0) {$y_1$};
                    \node[fact, label={[xshift=0.15cm,  yshift=0cm]left: $p(y_1 | x_1)$}] (fy1) at (0,1) {};
                    \node[cont] (y2) at (2,0) {$y_2$};
                    \node[fact, label={[xshift=0.15cm,  yshift=0cm]left: $p(y_2 | x_2)$}] (fy2) at (2,1) {};
                    \node[cont] (y3) at (4,0) {$y_3$};
                    \node[fact, label={[xshift=0.15cm,  yshift=0cm]left: $p(y_3 | x_3)$}] (fy3) at (4,1) {};
                    \node[cont] (y4) at (6,0) {$y_4$};
                    \node[fact, label={[xshift=0.15cm,  yshift=0cm]left: $p(y_4 | x_4)$}] (fy4) at (6,1) {};
                    
                    \node[fact, label=left: $p(x_1)$] (f1) at (-1,2) {};
                    \node[cont] (x1) at (0,2) {$x_1$};
                    \node[fact, label=above: $p(x_2 | x_1)$] (f2) at (1,2) {};
                    \node[cont] (x2) at (2,2) {$x_2$};
                    \node[fact, label=above: $p(x_3 | x_2)$] (f3) at (3,2) {};
                    \node[cont] (x3) at (4,2) {$x_3$};
                    \node[fact, label=above: $p(x_4 | x_3)$] (f4) at (5,2) {};
                    \node[cont] (x4) at (6,2) {$x_4$};
                    \draw(x1)--(y1);\draw(x2)--(y2);\draw(x3)--(y3);\draw(x4)--(y4);
                    \draw(x1)--(x2);\draw(x2)--(x3);\draw(x3)--(x4);
                    
                    \draw(f1) -- (x1);
                    
            \end{tikzpicture}}
        \end{center}
        
    \end{solution}
    
    \item Desenhe o grafo moralizado e um grafo de fatores n\~ao direcionado para modelos gr\'aficos direcionados definidos pelo grafo abaixo (este tipo de grafo \'e chamado de poli\'arvore: n\~ao h\'a loops, mas um n\'o pode ter mais de um pai).
    
    \begin{center}
        \scalebox{0.9}{ % x,y
            \begin{tikzpicture}[dgraph]
                \node[cont] (x1) at (0,0) {$x_1$};
                \node[cont] (x2) at (2,0) {$x_2$};
                \node[cont] (x3) at (-1,-1.5) {$x_3$};
                \node[cont] (x4) at (1,-1.5) {$x_4$};
                \node[cont] (x5) at (0,-3) {$x_5$};
                \node[cont] (x6) at (2,-3) {$x_6$};
                
                \draw(x1) -- (x3);
                \draw(x1) -- (x4);
                \draw(x2) -- (x4);
                \draw(x4) -- (x5);
                \draw(x4) -- (x6);
            \end{tikzpicture}
        }
    \end{center}
    
    \begin{solution}
        O grafo moral \'e obtido conectando os pais do n\'o colisor $x_4$. Veja o grafo \`a esquerda na figura abaixo.
        
        Para o grafo de fatores, notamos que o grafo direcionado define a
        seguinte classe de modelos probabil\'isticos:
        $$p(x_1, \ldots x_6) = p(x_1) p(x_2) p(x_3 | x_1) p(x_4|x_1,x_2) p(x_5|x_4) p(x_6|x_4)$$ 
        Isso resulta no grafo de fatores \`a direita na figura abaixo.
        
        \begin{center}
            \scalebox{0.9}{ % x,y
                \begin{tikzpicture}[ugraph]
                    \node[cont] (x1) at (0,0) {$x_1$};
                    \node[cont] (x2) at (2,0) {$x_2$};
                    \node[cont] (x3) at (-1,-1.5) {$x_3$};
                    \node[cont] (x4) at (1,-1.5) {$x_4$};
                    \node[cont] (x5) at (0,-3) {$x_5$};
                    \node[cont] (x6) at (2,-3) {$x_6$};
                    
                    \draw(x1) -- (x2);
                    \draw(x1) -- (x3);
                    \draw(x1) -- (x4);
                    \draw(x2) -- (x4);
                    \draw(x4) -- (x5);
                    \draw(x4) -- (x6);
                \end{tikzpicture}
                \hspace{10ex}
                \begin{tikzpicture}[ugraph]
                    \node[fact, label=left: $p(x_1)$] (f1) at (-1,0) {};      
                    \node[cont] (x1) at (0,0) {$x_1$};
                    
                    \node[fact, label=right: $p(x_2)$] (f2) at (3,0) {};      
                    \node[cont] (x2) at (2,0) {$x_2$};
                    
                    \node[fact, label=left: $p(x_3|x_1)$] (f3) at (-1,-0.7) {};      
                    \node[cont] (x3) at (-1,-1.5) {$x_3$};
                    
                    \node[fact, label=right: $p(x_4 | x_1\,x_2)$] (f4) at (1,-0.7) {};
                    \node[cont] (x4) at (1,-1.5) {$x_4$};
                    
                    \node[fact, label=left: $p(x_5 | x_4)$] (f5) at (0,-2.2) {};
                    \node[cont] (x5) at (0,-3) {$x_5$};
                    
                    \node[fact, label=right: $p(x_6 | x_4)$] (f6) at (2,-2.2) {};
                    \node[cont] (x6) at (2,-3) {$x_6$};
                    
                    \draw(f1) -- (x1);
                    \draw(f2) -- (x2);
                    \draw(x1) -- (f4);
                    \draw(x1) -- (f3);
                    
                    \draw(x2) -- (f4);
                    \draw(f4) -- (x4);
                    
                    \draw(f3) -- (x3);
                    
                    \draw(x4) -- (f5);
                    \draw(f5) -- (x5);
                    
                    \draw(x4) -- (f6);
                    \draw(f6) -- (x6);
                \end{tikzpicture}
            }
        \end{center}
        
        
        Nota:
        \begin{itemize}
            \item O grafo moral cont\'em um loop, enquanto o grafo de fatores n\~ao. O grafo de fatores ainda \'e uma poli\'arvore. Isso pode ser explorado para infer\^encia.
            \item Pode-se optar por agrupar alguns fatores para obter um grafo de fatores com uma estrutura específica (veja o grafo de fatores abaixo).
        \end{itemize}
        \vspace{1ex}
        
        \begin{center}
            \begin{tikzpicture}[ugraph]
                \node[cont] (x1) at (0,0) {$x_1$};
                
                \node[cont] (x2) at (2,0) {$x_2$};
                
                \node[fact, label=left: $p(x_3|x_1)$] (f3) at (-1,-0.7) {};      
                \node[cont] (x3) at (-1,-1.5) {$x_3$};
                
                \node[fact, label=right: $p(x_4 | x_1\,x_2)p(x_1)p(x_2)$] (f4) at (1,-0.7) {};
                \node[cont] (x4) at (1,-1.5) {$x_4$};
                
                \node[fact, label=right: $p(x_5 | x_4)p(x_6|x_4)$] (f5) at (1,-2.2) {};
                \node[cont] (x5) at (0,-3) {$x_5$};
                \node[cont] (x6) at (2,-3) {$x_6$};
                
                \draw(x1) -- (f4);
                \draw(x1) -- (f3);
                
                \draw(x2) -- (f4);
                \draw(f4) -- (x4);
                
                \draw(f3) -- (x3);
                
                \draw(x4) -- (f5);
                \draw(f5) -- (x5);
                \draw(f5) -- (x6);
            \end{tikzpicture}
        \end{center}
        
    \end{solution}
    
    
\end{exenumerate}



\ex{Passagem de mensagens sum-product}
\label{ex:sum-product-message-passing-alt}

Consideramos aqui a seguinte \'arvore de fatores:

\begin{center}
    \begin{tikzpicture}[ugraph]
        \node[fact, label=above: $\phi_A$] (fa) at (0,0) {} ; %
        \node[cont] (x1) at (1.5,0)  {$x_1$} ; %
        \node[fact, label=below: $\phi_C$] (fc) at (3,0) {} ; %
        \node[cont] (x2) at (3,1)  {$x_2$} ; %
        \node[fact, label=left: $\phi_B$] (fb) at (3,2) {} ; %
        \node[cont] (x3) at (4.5,0)  {$x_3$} ; %
        \node[fact, label=above: $\phi_D$] (fd) at (5.5,1) {} ; %
        \node[cont] (x4) at (7,1)  {$x_4$} ; %
        \node[fact, label=above: $\phi_E$] (fe) at (5.5,-1) {} ; %
        \node[cont] (x5) at (7,-1)  {$x_5$} ; %
        \node[fact, label=above: $\phi_F$] (ff) at (8.5,-1) {} ; %
        \draw (fa) -- (x1);
        \draw (x1) -- (fc);
        \draw (fc) -- (x2);
        \draw (x2) -- (fb);
        \draw (fc) -- (x3);
        \draw (x3) -- (fd);
        \draw (x3) -- (fe);
        \draw (fd) -- (x4);
        \draw (fe) -- (x5);
        \draw (x5) -- (ff);
    \end{tikzpicture}
\end{center}
Sejam todas as vari\'aveis bin\'arias, $x_i \in \{0,1\}$, e os fatores definidos como segue:
\begin{center}
    \begin{tabular}{ll}
        \toprule
        $x_1$ & $\phi_A$\\
        \midrule
        0 & 2\\
        1 & 4\\
        \bottomrule
    \end{tabular}
    \hfill
    \begin{tabular}{ll}
        \toprule
        $x_2$ & $\phi_B$\\
        \midrule
        0 & 4\\
        1 & 4\\
        \bottomrule
    \end{tabular}
    \hfill
    \begin{tabular}{llll}
        \toprule
        $x_1$ & $x_2$ & $x_3$ & $\phi_C$\\
        \midrule
        0 & 0 & 0 & 4 \\
        1 & 0 & 0 & 2 \\
        0 & 1 & 0 & 2 \\
        1 & 1 & 0 & 6 \\
        0 & 0 & 1 & 2 \\
        1 & 0 & 1 & 6 \\
        0 & 1 & 1 & 6 \\
        1 & 1 & 1 & 4 \\
        \bottomrule
    \end{tabular}
    \hfill
    \begin{tabular}{lll}
        \toprule
        $x_3$ & $x_4$ & $\phi_D$\\
        \midrule
        0 & 0 &  8 \\
        1 & 0 &  2 \\
        0 & 1 &  2 \\
        1 & 1 &  6 \\
        \bottomrule
    \end{tabular}
    \hfill
    \begin{tabular}{lll}
        \toprule
        $x_3$ & $x_5$ & $\phi_E$\\
        \midrule
        0 & 0 &  3 \\
        1 & 0 &  6 \\
        0 & 1 &  6 \\
        1 & 1 &  3 \\
        \bottomrule
    \end{tabular}
    \hfill
    \begin{tabular}{ll}
        \toprule
        $x_5$ & $\phi_F$\\
        \midrule
        0 & 1\\
        1 & 8\\
        \bottomrule
    \end{tabular}
    
    
\end{center}

\begin{exenumerate}
    
    \item Marque o grafo com setas indicando todas as mensagens que precisam
    ser computadas para o c\'alculo de $p(x_1)$.
    
    \begin{solution}
        
        \begin{center}
            \begin{tikzpicture}[ugraph]
                \node[fact, label=above: $\phi_A$] (fa) at (0,0) {} ; %
                \node[cont] (x1) at (1.5,0)  {$x_1$} ; %
                \node[fact, label=below: $\phi_C$] (fc) at (3,0) {} ; %
                \node[cont] (x2) at (3,1)  {$x_2$} ; %
                \node[fact, label=left: $\phi_B$] (fb) at (3,2) {} ; %
                \node[cont] (x3) at (4.5,0)  {$x_3$} ; %
                \node[fact, label=above: $\phi_D$] (fd) at (5.5,1) {} ; %
                \node[cont] (x4) at (7,1)  {$x_4$} ; %
                \node[fact, label=above: $\phi_E$] (fe) at (5.5,-1) {} ; %
                \node[cont] (x5) at (7,-1)  {$x_5$} ; %
                \node[fact, label=above: $\phi_F$] (ff) at (8.5,-1) {} ; %
                \draw (fa) -- (x1) node[midway,above,sloped] {$\rightarrow$};
                \draw (x1) -- (fc) node[midway,below,sloped] {$\leftarrow$};
                \draw (fc) -- (x2) node[midway,left] {$\downarrow$} ;
                \draw (x2) -- (fb) node[midway,left] {$\downarrow$};
                \draw (fc) -- (x3) node[midway,below] {$\leftarrow$};
                \draw (x3) -- (fd) node[midway,above,sloped] {$\leftarrow$};
                \draw (fd) -- (x4) node[midway,above] {$\leftarrow$};
                \draw (x3) -- (fe) node[midway,below,sloped] {$\leftarrow$};
                \draw (fe) -- (x5) node[midway,below] {$\leftarrow$};
                \draw (ff) -- (x5) node[midway,below] {$\leftarrow$};
            \end{tikzpicture}
        \end{center}
        
    \end{solution}
    
    \item \label{q:messages-alt} Compute as mensagens que voc\^e
    identificou.
    
    Assumindo que o c\'alculo das mensagens seja agendado de acordo
    com um rel\'ogio comum, agrupe as mensagens para que todas as mensagens
    no mesmo grupo possam ser computadas em paralelo durante um ciclo de
    rel\'ogio.
    
    \begin{solution}
        Como as vari\'aveis s\~ao bin\'arias, cada mensagem pode ser representada como um vetor bidimensional. Usamos a conven\'c\~ao de que o primeiro elemento do vetor corresponde \`a mensagem para $x_i = 0$ e o segundo elemento \`a mensagem para $x_i = 1$. Por exemplo,
        \begin{equation}
            \fxmessb{\phi_A}{x_1}= \begin{pmatrix}
                2 \\
                4 \\
            \end{pmatrix}
        \end{equation}
        significa que a mensagem $\fxmess{\phi_A}{x_1}{1}(x_1)$ \'e igual a 2 para $x_1=0$, i.e.\ $\fxmess{\phi_A}{x_1}{1}(0)=2$.
        
        A figura seguinte mostra um agrupamento (agendamento) do c\'alculo das mensagens. 
        \begin{figure}[h]
            \begin{center}
                \begin{tikzpicture}[ugraph]
                    \node[fact, label=above: $\phi_A$] (fa) at (0,0) {} ; %
                    \node[cont] (x1) at (1.5,0)  {$x_1$} ; %
                    \node[fact, label=below: $\phi_C$] (fc) at (3,0) {} ; %
                    \node[cont] (x2) at (3,1)  {$x_2$} ; %
                    \node[fact, label=left: $\phi_B$] (fb) at (3,2) {} ; %
                    \node[cont] (x3) at (4.5,0)  {$x_3$} ; %
                    \node[fact, label=above: $\phi_D$] (fd) at (5.5,1) {} ; %
                    \node[cont] (x4) at (7,1)  {$x_4$} ; %
                    \node[fact, label=above: $\phi_E$] (fe) at (5.5,-1) {} ; %
                    \node[cont] (x5) at (7,-1)  {$x_5$} ; %
                    \node[fact, label=above: $\phi_F$] (ff) at (8.5,-1) {} ; %
                    \draw (fa) -- (x1) node[midway,above,sloped] {$\substack{[1] \\ \rightarrow}$};
                    \draw (x1) -- (fc) node[midway,below,sloped] {$\substack{\leftarrow\\ [5]}$};
                    
                    \draw (fc) -- (x2) node[midway,left] {{\scriptsize $[2]$}$\downarrow$} ;
                    \draw (x2) -- (fb) node[midway,left] {{\scriptsize $[1]$}$\downarrow$};
                    \draw (fc) -- (x3) node[midway,below] {$\substack{\leftarrow \\ [4]}$};
                    \draw (x3) -- (fd) node[midway,above,sloped] {$\substack{[2] \\\leftarrow}$};
                    \draw (fd) -- (x4) node[midway,above] {$\substack{[1] \\ \leftarrow}$};
                    \draw (x3) -- (fe) node[midway,below,sloped] {$\substack{\leftarrow \\ [3]}$};
                    \draw (fe) -- (x5) node[midway,below] {$\substack{\leftarrow \\ [2]}$};
                    \draw (ff) -- (x5) node[midway,below] {$\substack{\leftarrow \\ [1]}$};
                \end{tikzpicture}
            \end{center}
        \end{figure}
        
        { \bf Ciclo de rel\'ogio 1:}\\
        \begin{align}
            \fxmessb{\phi_A}{x_1} & = \begin{pmatrix}
                2 \\
                4 \\
            \end{pmatrix}
            & 
            \fxmessb{\phi_B}{x_2} & = \begin{pmatrix}
                4 \\
                4 \\
            \end{pmatrix}
            &
            \xfmessb{x_4}{\phi_D} & = \begin{pmatrix}
                1 \\
                1 \\
            \end{pmatrix}
            &
            \fxmessb{\phi_F}{x_5} & = \begin{pmatrix}
                1 \\
                8 \\
            \end{pmatrix}
        \end{align}
        
        { \bf Ciclo de rel\'ogio 2:}\\
        \begin{align}
            \xfmessb{x_2}{\phi_C} = \fxmessb{\phi_B}{x_2}  & = \begin{pmatrix}
                4 \\
                4 \\
            \end{pmatrix}
            &
            \xfmessb{x_5}{\phi_E} =  \fxmessb{\phi_F}{x_5} & = \begin{pmatrix}
                1 \\
                8 \\
            \end{pmatrix}
        \end{align}
        A mensagem $\fxmess{\phi_D}{x_3}{1}$ \'e definida como
        \begin{align}
            \fxmess{\phi_D}{x_3}{1}(x_3) & = \sum_{x_4} \phi_D(x_3,x_4) \xfmess{x_4}{\phi_D}{1}(x_4)\end{align}
        de modo que
        \begin{align}
            \fxmess{\phi_D}{x_3}{1}(0) & = \sum_{x_4=0}^1 \phi_D(0,x_4) \xfmess{x_4}{\phi_D}{1}(x_4) \\
            & = \phi_D(0,0) \xfmess{x_4}{\phi_D}{1}(0) + \phi_D(0,1) \xfmess{x_4}{\phi_D}{1}(1) \\
            & = 8 \cdot 1  + 2 \cdot 1\\
            & = 10\\
            \fxmess{\phi_D}{x_3}{1}(1) & = \sum_{x_4=0}^1 \phi_D(1,x_4) \xfmess{x_4}{\phi_D}{1}(x_4) \\
            & = \phi_D(1,0) \xfmess{x_4}{\phi_D}{1}(0) + \phi_D(1,1) \xfmess{x_4}{\phi_D}{1}(1) \\
            & = 2 \cdot 1 + 6 \cdot 1\\
            & = 8
        \end{align}
        e assim
        \begin{align}
            \fxmessb{\phi_D}{x_3} &=  \begin{pmatrix}
                10 \\
                8
            \end{pmatrix}.
        \end{align}
        Os c\'alculos acima podem ser escritos de forma mais compacta em nota\'c\~ao de matriz. Seja $\pmb{\phi_D}$ a matriz que cont\'em os valores de $\phi_D(x_3, x_4)$ 
        \begin{align}
            \pmb{\phi_D} & = \begin{pmatrix}
                \phi_D(x_3=0,x_4=0) &  \phi_D(x_3=0,x_4=1) \\
                \phi_D(x_3=1,x_4=0) & \phi_D(x_3=1,x_4=1)  
            \end{pmatrix}
            = \begin{pmatrix}
                8 & 2 \\
                2 & 6
            \end{pmatrix}.
        \end{align}
        Podemos ent\~ao escrever $\fxmessb{\phi_D}{x_3}$ em termos de um produto matriz-vetor,
        \begin{align}
            \fxmessb{\phi_D}{x_3} & = \pmb{\phi_D} \xfmessb{x_4}{\phi_D}.
        \end{align}
        
        { \bf Ciclo de rel\'ogio 3:}\\
        Representando o fator $\phi_E$ como uma matriz $\pmb{\phi_E}$,
        \begin{align}
            \pmb{\phi_E} & = \begin{pmatrix}
                \phi_E(x_3=0,x_5=0) &  \phi_E(x_3=0,x_5=1) \\
                \phi_E(x_3=1,x_5=0) & \phi_E(x_3=1,x_5=1)  
            \end{pmatrix}
            = \begin{pmatrix}
                3 & 6 \\
                6 & 3
            \end{pmatrix},
        \end{align}
        podemos escrever
        \begin{align}
            \fxmess{\phi_E}{x_3}{1}(x_3) & = \sum_{x_5} \phi_E(x_3,x_5) \xfmess{x_5}{\phi_E}{1}(x_5)
        \end{align}
        como um produto matriz-vetor,
        \begin{align}
            \fxmessb{\phi_E}{x_3} & =  \pmb{\phi_E} \xfmessb{x_5}{\phi_E}\\
            & = \begin{pmatrix}
                3 & 6 \\
                6 & 3
            \end{pmatrix}
            \begin{pmatrix}
                1 \\
                8 \\
            \end{pmatrix}\\
            & =  \begin{pmatrix}
                51 \\
                30 \\
            \end{pmatrix}.
        \end{align}
        
        { \bf Ciclo de rel\'ogio 4:}\\
        O n\'o de vari\'avel $x_3$ recebeu todas as mensagens de entrada e pode, portanto, emitir $\xfmess{x_3}{\phi_C}{1}$,
        \begin{align}
            \xfmess{x_3}{\phi_C}{1}(x_3) & = \fxmess{\phi_D}{x_3}{1}(x_3)\fxmess{\phi_E}{x_3}{1}(x_3).
        \end{align}
        Usando $\odot$ para denotar a multiplica\'c\~ao elemento a elemento de dois vetores, temos 
        \begin{align}
            \xfmessb{x_3}{\phi_C} & = \fxmessb{\phi_D}{x_3} \odot \fxmessb{\phi_E}{x_3}\\
            & = \begin{pmatrix}
                10 \\
                8
            \end{pmatrix}
            \odot
            \begin{pmatrix}
                51 \\
                30 \\
            \end{pmatrix}\\
            & = \begin{pmatrix}
                510 \\
                240
            \end{pmatrix}.
        \end{align}
        
        { \bf Ciclo de rel\'ogio 5:}\\
        O n\'o de fator $\phi_C$ recebeu todas as mensagens de entrada e pode, portanto, emitir $\fxmess{\phi_C}{x_1}{1}$,
        \begin{align}
            \fxmess{\phi_C}{x_1}{1}(x_1) & = \sum_{x_2,x_3} \phi_C(x_1,x_2,x_3) \xfmess{x_2}{\phi_C}{1}(x_2)\xfmess{x_3}{\phi_C}{1}(x_3).
        \end{align}      
        Escrever a soma para $x_1=0$ e $x_1=1$ resulta em
        \begin{align}
            \fxmess{\phi_C}{x_1}{1}(0) = &  \sum_{x_2,x_3} \phi_C(0,x_2,x_3) \xfmess{x_2}{\phi_C}{1}(x_2) \xfmess{x_3}{\phi_C}{1}(x_3)\\
            = &  \phi_C(0,x_2,x_3) \xfmess{x_2}{\phi_C}{1}(x_2) \xfmess{x_3}{\phi_C}{1}(x_3) \mid_{(x_2,x_3)=(0,0)} +\\
            & \phi_C(0,x_2,x_3) \xfmess{x_2}{\phi_C}{1}(x_2) \xfmess{x_3}{\phi_C}{1}(x_3) \mid_{(x_2,x_3)=(1,0)} +\\
            & \phi_C(0,x_2,x_3) \xfmess{x_2}{\phi_C}{1}(x_2) \xfmess{x_3}{\phi_C}{1}(x_3) \mid_{(x_2,x_3)=(0,1)} +\\
            & \phi_C(0,x_2,x_3) \xfmess{x_2}{\phi_C}{1}(x_2) \xfmess{x_3}{\phi_C}{1}(x_3) \mid_{(x_2,x_3)=(1,1)}\\
            = & 4 \cdot 4 \cdot 510 +\\
            &  2 \cdot 4 \cdot 510 +  \\
            &  2 \cdot 4 \cdot 240  + \\
            &  6 \cdot 4 \cdot 240 \\
            = & 19920\\
            \fxmess{\phi_C}{x_1}{1}(1) = &  \sum_{x_2,x_3} \phi_C(1,x_2,x_3) \xfmess{x_2}{\phi_C}{1}(x_2) \xfmess{x_3}{\phi_C}{1}(x_3)\\
            = &  \phi_C(1,x_2,x_3) \xfmess{x_2}{\phi_C}{1}(x_2) \xfmess{x_3}{\phi_C}{1}(x_3) \mid_{(x_2,x_3)=(0,0)} +\\
            & \phi_C(1,x_2,x_3) \xfmess{x_2}{\phi_C}{1}(x_2) \xfmess{x_3}{\phi_C}{1}(x_3) \mid_{(x_2,x_3)=(1,0)} +\\
            & \phi_C(1,x_2,x_3) \xfmess{x_2}{\phi_C}{1}(x_2) \xfmess{x_3}{\phi_C}{1}(x_3) \mid_{(x_2,x_3)=(0,1)} +\\
            & \phi_C(1,x_2,x_3) \xfmess{x_2}{\phi_C}{1}(x_2) \xfmess{x_3}{\phi_C}{1}(x_3) \mid_{(x_2,x_3)=(1,1)}\\
            = & 2 \cdot 4 \cdot 510 +\\
            &  6 \cdot 4 \cdot 510 +  \\
            &  6 \cdot 4 \cdot 240  + \\
            &  4 \cdot 4 \cdot 240 \\
            = & 25920                 
        \end{align}
        e, portanto,
        \begin{align}
            \fxmessb{\phi_C}{x_1} &= \begin{pmatrix}
                19920\\
                25920
            \end{pmatrix}
        \end{align}
        
        Ap\'os o passo 5, o n\'o de vari\'avel $x_1$ recebeu todas as mensagens de entrada e a marginal pode ser calculada.
        
        Al\'em das mensagens necess\'arias para o c\'alculo de $p(x_1)$, pode-se
        calcular \emph{todas} as mensagens no grafo em cinco ciclos de rel\'ogio, veja
        a Figura \ref{fig:all_messages-alt}. Isso significa que \emph{todas} as marginais,
        bem como as conjuntas daquelas vari\'aveis que compartilham um n\'o de fator, est\~ao
        dispon\'iveis ap\'os cinco ciclos de rel\'ogio.
        
        \begin{figure}[h]
            \begin{center}
                \scalebox{1.6}{
                    \begin{tikzpicture}[ugraph]
                        \node[fact, label=above: {\tiny $\phi_A$}] (fa) at (0,0) {} ; %
                        \node[cont] (x1) at (1.5,0)  {$x_1$} ; %
                        \node[fact, label=below:  {\tiny$\phi_C$}] (fc) at (3,0) {} ; %
                        \node[cont] (x2) at (3,1)  {$x_2$} ; %
                        \node[fact, label=left: {\tiny $\phi_B$}] (fb) at (3,2) {} ; %
                        \node[cont] (x3) at (4.5,0)  {$x_3$} ; %
                        \node[fact, label=above: {\tiny  $\phi_D$}] (fd) at (5.5,1) {} ; %
                        \node[cont] (x4) at (7,1)  {$x_4$} ; %
                        \node[fact, label=above: {\tiny  $\phi_E$}] (fe) at (5.5,-1) {} ; %
                        \node[cont] (x5) at (7,-1)  {$x_5$} ; %
                        \node[fact, label=above:  {\tiny $\phi_F$}] (ff) at (8.5,-1) {} ; %
                        \draw (fa) -- (x1) node[midway,above,sloped] {$\substack{[1] \\ \rightarrow}$};
                        \draw (x1) -- (fc) node[midway,below,sloped] {$\substack{\leftarrow\\ [5]}$} node[near start,above,sloped] {$\substack{[2] \\ \rightarrow}$};
                        
                        \draw (fc) -- (x2) node[midway,left] {{\tiny $[2]$}$\downarrow$}  node[midway,right] {$\uparrow${\tiny $[5]$}} ;
                        \draw (x2) -- (fb) node[midway,left] {{\tiny $[1]$}$\downarrow$};
                        \draw (fc) -- (x3) node[midway,below] {$\substack{\leftarrow \\ [4]}$} node[near end,above] {$\substack{[3] \\ \rightarrow }$};
                        \draw (x3) -- (fd) node[midway,above,sloped] {$\substack{[2] \\\leftarrow}$} node[midway,below,sloped] {$\substack{\rightarrow \\ [4]}$};
                        \draw (fd) -- (x4) node[midway,above] {$\substack{[1] \\ \leftarrow}$} node[midway,below] {$\substack{\rightarrow \\ [5]}$}; 
                        \draw (x3) -- (fe) node[midway,below,sloped] {$\substack{\leftarrow \\ [3]}$}  node[midway,above,sloped] {$\substack{[4] \\ \rightarrow}$};
                        \draw (fe) -- (x5) node[midway,below] {$\substack{\leftarrow \\ [2]}$} node[midway,above] {$\substack{ [5] \\ \rightarrow }$}; 
                        \draw (ff) -- (x5) node[midway,below] {$\substack{\leftarrow \\ [1]}$};
                \end{tikzpicture}}
            \end{center}
            \caption{\label{fig:all_messages-alt} Resposta ao Exerc\'icio \ref{ex:sum-product-message-passing-alt} Quest\~ao \ref{q:messages-alt}: Calculando todas as mensagens em cinco ciclos de rel\'ogio. Se tamb\'em comput\'assemos as mensagens em dire\'c\~ao aos n\'os de fator folha, seriam necess\'arios seis ciclos, mas elas n\~ao s\~ao necess\'arias para o c\'alculo das marginais, por isso foram omitidas.}
        \end{figure}
        
        
    \end{solution}
    
    \item Qual \'e $p(x_1=1)$?
    
    \begin{solution}
        Calculamos a marginal $p(x_1)$ como
        \begin{align}
            p(x_1) \propto \fxmess{\phi_A}{x_1}{1}(x_1) \fxmess{\phi_C}{x_1}{1}(x_1)
        \end{align}
        que \'e, em nota\'c\~ao de vetor,
        \begin{align}
            \begin{pmatrix}
                p(x_1=0)\\
                p(x_1=1)
            \end{pmatrix}
            & \propto 
            \fxmessb{\phi_A}{x_1} \odot \fxmessb{\phi_C}{x_1}\\
            & \propto  
            \begin{pmatrix}
                2 \\
                4 \\
            \end{pmatrix}
            \odot
            \begin{pmatrix}
                19920\\
                25920
            \end{pmatrix}\\
            & \propto
            \begin{pmatrix}
                39840\\
                103680
            \end{pmatrix}.
        \end{align}
        A normaliza\'c\~ao resulta em
        \begin{align}
            \begin{pmatrix}
                p(x_1=0)\\
                p(x_1=1)
            \end{pmatrix}
            & =  \frac{1}{39840+103680}\begin{pmatrix}
                39840\\
                103680
            \end{pmatrix}\\
            & = \begin{pmatrix}
                0.2776\\
                0.7224
            \end{pmatrix}
        \end{align}
        de modo que $p(x_1=1) = 0.7224$.
        
        Note os n\'umeros relativamente grandes nas mensagens que
        computamos. Em outros casos, pode-se obter valores muito pequenos,
        dependendo da escala dos fatores. Isso pode causar problemas
        num\'ericos que podem ser resolvidos trabalhando no dom\'inio
        logar\'itmico.
    \end{solution}
    
    \item Desenhe o grafo de fatores correspondente a $p(x_1, x_3, x_4, x_5 | x_2 =1)$ e forne\'ca os valores num\'ericos para todos os fatores.
    
    \begin{solution}
        A fmp representada pelo grafo de fatores original \'e
        $$p(x_1, \ldots, x_5) \propto \phi_A(x_1) \phi_B(x_2) \phi_C(x_1,x_2,x_3) \phi_D(x_3,x_4) \phi_E(x_3,x_5) \phi_F(x_5)$$
        A condicional $p(x_1, x_3, x_4, x_5 | x_2 =1)$ \'e proporcional a $p(x_1, \ldots, x_5)$ com $x_2$ fixado em $x_2=1$, i.e.\
        \begin{align}
            p(x_1, x_3, x_4, x_5 | x_2 =1) &\propto p(x_1,x_2=1, x_3, x_4, x_5) \\
            & \propto \phi_A(x_1) \phi_B(x_2=1) \phi_C(x_1,x_2=1,x_3) \phi_D(x_3,x_4) \phi_E(x_3,x_5) \phi_F(x_5)\\
            & \propto \phi_A(x_1) \phi^{x_2}_C(x_1,x_3) \phi_D(x_3,x_4) \phi_E(x_3,x_5) \phi_F(x_5)
        \end{align}
        onde $\phi^{x_2}_C(x_1,x_3) = \phi_C(x_1,x_2=1,x_3)$. Os valores num\'ericos de $\phi^{x_2}_C(x_1,x_3)$ podem ser lidos da tabela que define $\phi_C(x_1,x_2,x_3)$, extraindo aquelas linhas onde $x_2=1$,
        \begin{center}
            \begin{tabular}{lllll}
                \toprule
                & $x_1$ & $x_2$ & $x_3$ & $\phi_C$\\
                \midrule
                &0 & 0 & 0 & 4 \\
                &1 & 0 & 0 & 2 \\
                $\rightarrow$&0 & 1 & 0 & 2 \\
                $\rightarrow$& 1 & 1 & 0 & 6 \\
                &0 & 0 & 1 & 2 \\
                &1 & 0 & 1 & 6 \\
                $\rightarrow$&0 & 1 & 1 & 6 \\
                $\rightarrow$&1 & 1 & 1 & 4 \\
                \bottomrule
            \end{tabular}
            \hspace{3ex} \text{de modo que} \hspace{3ex}
            \begin{tabular}{lll}
                \toprule
                $x_1$ & $x_3$ & $\phi^{x_2}_C$\\
                \midrule
                0 & 0 & 2 \\
                1 & 0 & 6 \\
                0 & 1 & 6 \\
                1 & 1 & 4 \\
                \bottomrule
            \end{tabular}
        \end{center}
        
        O grafo de fatores para $p(x_1, x_3, x_4, x_5 | x_2 =1)$ \'e mostrado abaixo. O fator $\phi_B$ desapareceu, pois dependia apenas de $x_2$ e, portanto, tornou-se uma constante. O fator $\phi_C$ \'e substitu\'ido por $\phi_C^{x_2}$ definido acima. Os fatores restantes s\~ao os mesmos do grafo de fatores original.
        
        \begin{center}
            \begin{tikzpicture}[ugraph]
                \node[fact, label=above: $\phi_A$] (fa) at (0,0) {} ; %
                \node[cont] (x1) at (1.5,0)  {$x_1$} ; %
                \node[fact, label=below: $\phi^{x_2}_C$] (fc) at (3,0) {} ; %
                %\node[cont] (x2) at (3,1)  {$x_2$} ; %
                %\node[fact, label=left: $\phi_B$] (fb) at (3,2) {} ; %
                \node[cont] (x3) at (4.5,0)  {$x_3$} ; %
                \node[fact, label=above: $\phi_D$] (fd) at (5.5,1) {} ; %
                \node[cont] (x4) at (7,1)  {$x_4$} ; %
                \node[fact, label=above: $\phi_E$] (fe) at (5.5,-1) {} ; %
                \node[cont] (x5) at (7,-1)  {$x_5$} ; %
                \node[fact, label=above: $\phi_F$] (ff) at (8.5,-1) {} ; %
                \draw (fa) -- (x1);
                \draw (x1) -- (fc);
                %\draw (fc) -- (x2);
                %\draw (x2) -- (fb);
                \draw (fc) -- (x3);
                \draw (x3) -- (fd);
                \draw (x3) -- (fe);
                \draw (fd) -- (x4);
                \draw (fe) -- (x5);
                \draw (x5) -- (ff);
            \end{tikzpicture}
        \end{center}
        
    \end{solution}
    
    \item Calcule $p(x_1 =1 | x_2 = 1)$, reutilizando mensagens que voc\^e j\'a calculou para a avalia\'c\~ao de $p(x_1 =1)$.
    
    \begin{solution}
        A mensagem $\fxmess{\phi_A}{x_1}{1}$ \'e a mesma do grafo de fatores
        original e $\xfmess{x_3}{\phi_C^{x_2}}{1} = \xfmess{x_3}{\phi_C}{1}$. Isso
        ocorre porque a mensagem de sa\'ida de $x_3$ corresponde ao fator
        efetivo obtido ao somar todas as vari\'aveis nas sub\'arvores anexadas a
        $x_3$ (sem o ramo $\phi_C^{x_2}$), e essas sub\'arvores n\~ao dependem de $x_2$.
        
        A mensagem $\fxmess{\phi_C^{x_2}}{x_1}{1}$ precisa ser calculada novamente. Temos
        \begin{align}
            \fxmess{\phi_C^{x_2}}{x_1}{1}(x_1) & = \sum_{x_3} \phi_C^{x_2}(x_1,x_3) \xfmess{x_3}{\phi_C^{x_2}}{1}
        \end{align}
        ou em nota\'c\~ao de vetor
        \begin{align}
            \fxmessb{\phi_C^{x_2}}{x_1} & = \pmb{\phi_C^{x_2}} \xfmessb{x_3}{\phi_C^{x_2}}\\
            &= \begin{pmatrix}
                \phi_C^{x_2}(x_1=0,x_3=0) &  \phi_C^{x_2}(x_1=0,x_3=1) \\
                \phi_C^{x_2}(x_1=1,x_3=0) & \phi_C^{x_2}(x_1=1,x_3=1)  
            \end{pmatrix}
            \xfmessb{x_3}{\phi_C^{x_2}}\\
            & = \begin{pmatrix}
                2 & 6 \\
                6 & 4
            \end{pmatrix}
            \begin{pmatrix}
                510\\
                240
            \end{pmatrix}\\
            &=
            \begin{pmatrix}
                2460\\
                4020
            \end{pmatrix}
        \end{align}
        Obtemos assim a posterior marginal de $x_1$ dado $x_2=1$:
        \begin{align}
            \begin{pmatrix}
                p(x_1=0 | x_2=1)\\
                p(x_1=1 | x_2=1)
            \end{pmatrix}
            & \propto 
            \fxmessb{\phi_A}{x_1} \odot \fxmessb{\phi^{x_2}_C}{x_1}\\
            & \propto  
            \begin{pmatrix}
                2 \\
                4 \\
            \end{pmatrix}
            \odot
            \begin{pmatrix}
                2460\\
                4020
            \end{pmatrix}\\
            & \propto
            \begin{pmatrix}
                4920\\
                16080
            \end{pmatrix}.
        \end{align}
        A normaliza\'c\~ao resulta em
        \begin{align}
            \begin{pmatrix}
                p(x_1=0 | x_2=1)\\
                p(x_1=1 | x_2=1)
            \end{pmatrix}
            & =  \begin{pmatrix}
                0.2343\\
                0.7657
            \end{pmatrix}
        \end{align}
        e, portanto, $p(x_1=1 | x_2 =1) = 0.7657$. A probabilidade posterior \'e ligeiramente maior que a probabilidade a priori, $p(x_1=1) = 0.7224$.
    \end{solution}
    
\end{exenumerate}



\ex{Passagem de mensagens sum-product}
O seguinte grafo de fatores representa uma distribui\'c\~ao de Gibbs
sobre quatro vari\'aveis bin\'arias $x_i \in \{0,1\}$.
\begin{figure}[h]
    \begin{center}
        \scalebox{1}{ % x,y
            \begin{tikzpicture}[ugraph]
                \node[fact, label=above: $\phi_a$] (fa) at (-1.5,0) {};
                \node[cont] (x1) at (0,0) {$x_1$};
                
                \node[fact, label=above: $\phi_b$] (fb) at (0,1.5) {};
                \node[cont] (x2) at (-1.5,1.5) {$x_2$};
                
                \node[fact, label=above: $\phi_c$] (fc) at (1.5,0) {};
                \node[cont] (x3) at (3,1.5) {$x_3$};
                \node[cont] (x4) at (3,-1.5) {$x_4$};
                
                \node[fact, label=above: $\phi_d$] (fd) at (4.5,1.5) {};
                \node[fact, label=above: $\phi_e$] (fe) at (4.5,-1.5) {};
                
                \draw (fa) -- (x1);
                \draw (fb) -- (x1);
                \draw (fb) -- (x2);
                \draw (x1) -- (fc);
                \draw (fc) -- (x3);
                \draw (fc) -- (x4);
                \draw (x3) -- (fd);
                \draw (x4) -- (fe);
        \end{tikzpicture}}
    \end{center}
\end{figure}\\
Os fatores $\phi_a, \phi_b, \phi_d$ s\~ao definidos como segue:\\
\begin{center}
    \begin{tabular}{l l}
        \toprule
        $x_1$ & $\phi_a$\\
        \midrule
        0 & 2 \\
        1 & 1\\
        \bottomrule
    \end{tabular}
    \hspace{10ex}
    \begin{tabular}{l l l}
        \toprule
        $x_1$ & $x_2$ & $\phi_b$\\
        \midrule
        0 & 0 & 5 \\
        1 & 0 & 2 \\
        0 & 1 & 2 \\
        1 & 1 & 6 \\
        \bottomrule
    \end{tabular}
    \hspace{10ex}
    \begin{tabular}{l l}
        \toprule
        $x_3$ & $\phi_d$\\
        \midrule
        0 & 1 \\
        1 & 2\\
        \bottomrule
    \end{tabular}\\
\end{center}
e $\phi_c(x_1,x_3,x_4) = 1$ se $x_1=x_3=x_4$, e \'e zero caso contr\'ario.\\[2ex]
Para todas as quest\~oes abaixo, justifique sua resposta:
\begin{exenumerate}
    \item Calcule os valores de $\xfmess{x_2}{\phi_b}{1}(x_2)$ para $x_2=0$ e $x_2=1$.
    
    \begin{solution}
        Mensagens de n\'os de vari\'avel folha para n\'os de fator s\~ao iguais a um, de modo que $\xfmess{x_2}{\phi_b}{1}(x_2)=1$ para todo $x_2$.
    \end{solution}
    
    \item Assuma que a mensagem $\xfmess{x_4}{\phi_c}{1}(x_4)$ seja igual a
    $$\xfmess{x_4}{\phi_c}{1}(x_4) = \begin{cases}
        1 & \text{se } x_4=0\\
        3 & \text{se } x_4=1\\
    \end{cases}$$
    Calcule os valores de $\phi_e(x_4)$ para $x_4=0$ e $x_4=1$. 
    
    \begin{solution}
        
        Mensagens de fatores folha para seus n\'os de vari\'avel s\~ao
        iguais aos fatores folha, e n\'os de vari\'avel com uma \'unica
        mensagem de entrada copiam a mensagem. Temos, portanto,
        \begin{align}
            \fxmess{\phi_e}{x_4}{1}(x_4) & = \phi_e(x_4)\\
            \xfmess{x_4}{\phi_c}{1}(x_4) & =  \fxmess{\phi_e}{x_4}{1}(x_4)
        \end{align}
        e assim
        \begin{align}
            \phi_e(x_4) &= \begin{cases}
                1 & \text{se } x_4=0\\
                3 & \text{se } x_4=1\\
            \end{cases}
        \end{align}
        
    \end{solution}
    
    
    \item Calcule os valores de $\fxmess{\phi_c}{x_1}{1}(x_1)$ para $x_1=0$ e $x_1=1$.
    
    \begin{solution}
        
        Primeiro calculamos $\xfmess{x_3}{\phi_c}{1}(x_3)$:
        \begin{align}
            \xfmess{x_3}{\phi_c}{1}(x_3) & = \fxmess{\phi_d}{x_3}{1}(x_3)\\
            & = \begin{cases}
                1 & \text{se } x_3=0\\
                2 & \text{se } x_3=1
            \end{cases}
        \end{align}
        A mensagem desejada $\fxmess{\phi_c}{x_1}{1}(x_1)$ \'e, por defini\'c\~ao,
        \begin{align}
            \fxmess{\phi_c}{x_1}{1}(x_1) & = \sum_{x_3,x_4} \phi_c(x_1,x_3,x_4) \xfmess{x_3}{\phi_c}{1}(x_3) \xfmess{x_4}{\phi_c}{1}(x_4)
        \end{align}
        Como $\phi_c(x_1,x_3,x_4)$ s\'o \'e n\~ao nulo se $x_1=x_3=x_4$, onde \'e igual a um, os c\'alculos simplificam-se:
        \begin{align}
            \fxmess{\phi_c}{x_1}{1}(x_1=0) & = \phi_c(0,0,0) \xfmess{x_3}{\phi_c}{1}(0) \xfmess{x_4}{\phi_c}{1}(0)\\
            & = 1 \cdot 1 \cdot 1\\
            & = 1
        \end{align}
        \begin{align}
            \fxmess{\phi_c}{x_1}{1}(x_1=1) & = \phi_c(1,1,1) \xfmess{x_3}{\phi_c}{1}(1) \xfmess{x_4}{\phi_c}{1}(1)\\
            & = 1 \cdot 2 \cdot 3\\
            & = 6
        \end{align}
    \end{solution}
    
    
    \item A mensagem $\fxmess{\phi_b}{x_1}{1}(x_1)$ \'e igual a
    $$ \fxmess{\phi_b}{x_1}{1}(x_1) = \begin{cases}
        7 & \text{se } x_1=0\\
        8 & \text{se } x_1=1\\
    \end{cases} $$
    Qual \'e a probabilidade de que $x_1=1$, i.e.\ $p(x_1=1)$?
    
    \begin{solution}
        A marginal n\~ao normalizada $p(x_1)$ \'e dada pelo produto das tr\^es mensagens de entrada
        \begin{align}
            p(x_1) & \propto \fxmess{\phi_a}{x_1}{1}(x_1)\fxmess{\phi_b}{x_1}{1}(x_1) \fxmess{\phi_c}{x_1}{1}(x_1)
        \end{align}
        Com
        \begin{align}
            \fxmess{\phi_b}{x_1}{1}(x_1) & = \sum_{x_2} \phi_b(x_1,x_2)
        \end{align}
        segue que
        \begin{align}
            \fxmess{\phi_b}{x_1}{1}(x_1=0) & = \sum_{x_2} \phi_b(0,x_2)\\
            & = 5+2 \\
            & = 7\\
            \fxmess{\phi_b}{x_1}{1}(x_1=1) & = \sum_{x_2} \phi_b(1,x_2)\\
            & = 2+6\\
            & = 8
        \end{align}
        Portanto, obtemos
        \begin{align}
            p(x_1=0) &\propto 2 \cdot 7 \cdot 1 = 14\\
            p(x_1=1) &\propto 1 \cdot 8 \cdot 6 = 48
        \end{align}
        e a normaliza\'c\~ao fornece o resultado desejado
        \begin{align}
            p(x_1=1) & = \frac{48}{14+48} = \frac{48}{62} = \frac{24}{31} \approx 0.774
        \end{align}
        
    \end{solution}
    
    
    
\end{exenumerate}




\ex{Passagem de mensagens max-sum}
\label{ex:max-sum-message-passing-alt}
Calculamos aqui os estados mais prov\'aveis para o grafo de fatores e os fatores abaixo.

\begin{center}
    \begin{tikzpicture}[ugraph]
        \node[fact, label=above: $\phi_A$] (fa) at (0,0) {} ; %
        \node[cont] (x1) at (1.5,0)  {$x_1$} ; %
        \node[fact, label=below: $\phi_C$] (fc) at (3,0) {} ; %
        \node[cont] (x2) at (3,1)  {$x_2$} ; %
        \node[fact, label=left: $\phi_B$] (fb) at (3,2) {} ; %
        \node[cont] (x3) at (4.5,0)  {$x_3$} ; %
        \node[fact, label=above: $\phi_D$] (fd) at (5.5,1) {} ; %
        \node[cont] (x4) at (7,1)  {$x_4$} ; %
        \node[fact, label=above: $\phi_E$] (fe) at (5.5,-1) {} ; %
        \node[cont] (x5) at (7,-1)  {$x_5$} ; %
        \node[fact, label=above: $\phi_F$] (ff) at (8.5,-1) {} ; %
        \draw (fa) -- (x1);
        \draw (x1) -- (fc);
        \draw (fc) -- (x2);
        \draw (x2) -- (fb);
        \draw (fc) -- (x3);
        \draw (x3) -- (fd);
        \draw (x3) -- (fe);
        \draw (fd) -- (x4);
        \draw (fe) -- (x5);
        \draw (x5) -- (ff);
    \end{tikzpicture}
\end{center}

Sejam todas as vari\'aveis bin\'arias, $x_i \in \{0,1\}$, e os fatores definidos como segue:
\begin{center}
    \begin{tabular}{ll}
        \toprule
        $x_1$ & $\phi_A$\\
        \midrule
        0 & 2\\
        1 & 4\\
        \bottomrule
    \end{tabular}
    \hfill
    \begin{tabular}{ll}
        \toprule
        $x_2$ & $\phi_B$\\
        \midrule
        0 & 4\\
        1 & 4\\
        \bottomrule
    \end{tabular}
    \hfill
    \begin{tabular}{llll}
        \toprule
        $x_1$ & $x_2$ & $x_3$ & $\phi_C$\\
        \midrule
        0 & 0 & 0 & 4 \\
        1 & 0 & 0 & 2 \\
        0 & 1 & 0 & 2 \\
        1 & 1 & 0 & 6 \\
        0 & 0 & 1 & 2 \\
        1 & 0 & 1 & 6 \\
        0 & 1 & 1 & 6 \\
        1 & 1 & 1 & 4 \\
        \bottomrule
    \end{tabular}
    \hfill
    \begin{tabular}{lll}
        \toprule
        $x_3$ & $x_4$ & $\phi_D$\\
        \midrule
        0 & 0 &  8 \\
        1 & 0 &  2 \\
        0 & 1 &  2 \\
        1 & 1 &  6 \\
        \bottomrule
    \end{tabular}
    \hfill
    \begin{tabular}{lll}
        \toprule
        $x_3$ & $x_5$ & $\phi_E$\\
        \midrule
        0 & 0 &  3 \\
        1 & 0 &  6 \\
        0 & 1 &  6 \\
        1 & 1 &  3 \\
        \bottomrule
    \end{tabular}
    \hfill
    \begin{tabular}{ll}
        \toprule
        $x_5$ & $\phi_F$\\
        \midrule
        0 & 1\\
        1 & 8\\
        \bottomrule
    \end{tabular}
    
\end{center}

\begin{exenumerate}
    \item Precisaremos calcular a constante de normaliza\'c\~ao $Z$ para determinar $\argmax_{\x} p(x_1, \ldots, x_5)$?
    
    \begin{solution}
        Isso n\~ao \'e necess\'ario, j\'a que $\argmax_{\x} p(x_1, \ldots, x_5) = \argmax_{\x} c p(x_1, \ldots, x_5)$ para qualquer constante $c$. Algoritmicamente, o algoritmo de backtracking tamb\'em \'e invariante a qualquer escala dos fatores.
    \end{solution}
    \item Calcule $\argmax_{x_1, x_2, x_3} p(x_1, x_2, x_3 | x_4=0, x_5=0)$ via passagem de mensagens max-sum.
    \begin{solution}
        Primeiro derivamos o grafo de fatores e os fatores correspondentes para
        $p(x_1, x_2, x_3 | x_4=0, x_5=0)$.
        
        Para valores fixos de $x_4, x_5$, as duas vari\'aveis s\~ao removidas
        do grafo, e os fatores $\phi_D(x_3,x_4)$ e $\phi_E(x_3, x_5)$ s\~ao
        reduzidos a fatores univariados $\phi_D^{x_4}(x_3)$ e $\phi_D^{x_5}(x_3)$
        ao reter aquelas linhas na tabela onde $x_4=0$ e $x_5=0$, respectivamente:
        \begin{center}
            \begin{tabular}{ll}
                \toprule
                $x_3$ & $\phi_D^{x_4}$\\
                \midrule
                0 & 8\\
                1 & 2\\
                \bottomrule
            \end{tabular}
            \hspace{4ex}
            \begin{tabular}{ll}
                \toprule
                $x_3$ & $\phi_E^{x_5}$\\
                \midrule
                0 & 3\\
                1 & 6\\
                \bottomrule
            \end{tabular}   
        \end{center}
        Como ambos os fatores dependem apenas de $x_3$, eles podem ser combinados em um novo fator $\tilde{\phi}(x_3)$ por multiplica\'c\~ao elemento a elemento.
        \begin{center}
            \begin{tabular}{ll}
                \toprule
                $x_3$ & $\tilde{\phi}$\\
                \midrule
                0 & 24\\
                1 & 12\\
                \bottomrule
            \end{tabular}
        \end{center}
        Al\'em disso, como trabalhamos com um modelo n\~ao normalizado, podemos reescalar o fator para que o valor m\'aximo seja um (ou qualquer outro valor conveniente), de modo que:
        \begin{center}
            \begin{tabular}{ll}
                \toprule
                $x_3$ & $\tilde{\phi}$\\
                \midrule
                0 & 2\\
                1 & 1\\
                \bottomrule
            \end{tabular}
        \end{center}
        O fator $\phi_F(x_5)$ \'e uma constante para um valor fixo de $x_5$ e pode ser ignorado. O grafo de fatores para $p(x_1, x_2, x_3 | x_4=0, x_5=0)$ \'e, portanto:
        \begin{center}
            \begin{tikzpicture}[ugraph]
                \node[fact, label=above: $\phi_A$] (fa) at (0,0) {} ; %
                \node[cont] (x1) at (1.5,0)  {$x_1$} ; %
                \node[fact, label=below: $\phi_C$] (fc) at (3,0) {} ; %
                \node[cont] (x2) at (3,1)  {$x_2$} ; %
                \node[fact, label=left: $\phi_B$] (fb) at (3,2) {} ; %
                \node[cont] (x3) at (4.5,0)  {$x_3$} ; %
                \node[fact, label=above: $\tilde{\phi}$] (fnew) at (6,0) {} ; %
                \draw (fa) -- (x1);
                \draw (x1) -- (fc);
                \draw (fc) -- (x2);
                \draw (x2) -- (fb);
                \draw (fc) -- (x3);
                \draw (x3) -- (fnew);
            \end{tikzpicture}
        \end{center}
        
        Vamos fixar $x_1$ como raiz em dire\'c\~ao \`a qual as mensagens s\~ao computadas. As mensagens que precisamos computar s\~ao mostradas no seguinte grafo:
        \begin{center}
            \begin{tikzpicture}[ugraph]
                \node[fact, label=above: $\phi_A$] (fa) at (0,0) {} ; %
                \node[cont] (x1) at (1.5,0)  {$x_1$} ; %
                \node[fact, label=below: $\phi_C$] (fc) at (3,0) {} ; %
                \node[cont] (x2) at (3,1)  {$x_2$} ; %
                \node[fact, label=left: $\phi_B$] (fb) at (3,2) {} ; %
                \node[cont] (x3) at (4.5,0)  {$x_3$} ; %
                \node[fact, label=above: $\tilde{\phi}$] (fnew) at (6,0) {} ; %
                \draw (fa) -- (x1) node[midway,above] {$\rightarrow$}; 
                \draw (x1) -- (fc) node[midway,above] {$\leftarrow$}; 
                \draw (fc) -- (x2) node[midway,left] {$\downarrow$}; 
                \draw (x2) -- (fb) node[midway,left] {$\downarrow$}; 
                \draw (fc) -- (x3) node[midway,above] {$\leftarrow$}; 
                \draw (x3) -- (fnew) node[midway,above] {$\leftarrow$}; 
            \end{tikzpicture}
        \end{center}
        
        Em seguida, calculamos as mensagens das folhas (log-messages). Temos apenas n\'os de fator como n\'os folha, de modo que:
        \begin{align}
            \lambdab_{\phi_A \to x_1} & =  \begin{pmatrix}
                \log \phi_A(x_1=0)\\
                \log \phi_A(x_1=1)
            \end{pmatrix}
            = \begin{pmatrix}
                \log 2\\
                \log 4
            \end{pmatrix}
        \end{align}
        e similarmente:
        \begin{align}
            \lambdab_{\phi_B \to x_2} & =   \begin{pmatrix}
                \log \phi_B(x_2=0)\\
                \log \phi_B(x_2=1)
            \end{pmatrix} = \begin{pmatrix}
                \log 4\\
                \log 4
            \end{pmatrix}
            &
            \lambdab_{\tilde{\phi} \to x_3} & =   \begin{pmatrix}
                \log \tilde{\phi}(x_3=0)\\
                \log \tilde{\phi}(x_3=1)
            \end{pmatrix}= \begin{pmatrix}
                \log 2\\
                \log 1
            \end{pmatrix}
        \end{align}
        Como os n\'os de vari\'avel $x_2$ e $x_3$ t\^em apenas uma aresta de entrada cada, obtemos:
        \begin{align}
            \lambdab_{x_2 \to \phi_C} = \lambdab_{\phi_B \to x_2} & =  \begin{pmatrix}
                \log 4\\
                \log 4
            \end{pmatrix}
            &
            \lambdab_{x_3 \to \phi_C} = \lambdab_{\tilde{\phi} \to x_3} & =  \begin{pmatrix}
                \log 2\\
                \log 1
            \end{pmatrix}
        \end{align}
        A mensagem $\lambda_{\phi_C \to x_1}(x_1)$ \'e igual a:
        \begin{align}
            \lambda_{\phi_C \to x_1}(x_1) & = \max_{x_2, x_3} \{ \log \phi_C(x_1, x_2, x_3) + \lambda_{x_2 \to \phi_C}(x_2) + \lambda_{x_3 \to \phi_C}(x_3) \}
        \end{align}
        onde escrevemos as mensagens em nota\'c\~ao n\~ao vetorial para destacar sua depend\^encia das vari\'aveis $x_2$ e $x_3$. Agora temos que considerar todas as combina\'c\~oes de $x_2$ e $x_3$:
        \begin{center}
            \begin{tabular}{lll}
                \toprule
                $x_2$ & $x_3$ & $\log \phi_C(x_1=0, x_2, x_3)$\\
                \midrule
                0 & 0 & $\log 4$ \\
                1 & 0 &  $\log 2$\\
                0 & 1 & $\log 2$ \\
                1 & 1 & $\log 6$ \\
                \bottomrule
            \end{tabular}
            \hspace{4ex}
            \begin{tabular}{lll}
                \toprule
                $x_2$ & $x_3$ & $\log \phi_C(x_1=1, x_2, x_3)$\\
                \midrule
                0 & 0 & $\log 2$  \\
                1 & 0 & $\log 6 $ \\
                0 & 1 & $\log 6$  \\
                1 & 1 & $\log 4$ \\
                \bottomrule
            \end{tabular}
        \end{center}
        Al\'em disso:
        \begin{center}
            \begin{tabular}{lll}
                \toprule
                $x_2$ & $x_3$ & $\lambda_{x_2 \to \phi_C}(x_2) + \lambda_{x_3 \to \phi_C}(x_3)$\\
                \midrule
                0 & 0 & $\log 4 + \log 2 = \log 8$  \\
                1 & 0 & $\log 4 + \log 2 = \log 8$\\
                0 & 1 & $\log 4 + 0 = \log 4$\\
                1 & 1 & $\log 4 + 0 = \log 4$\\
                \bottomrule
            \end{tabular}
        \end{center}
        Portanto, para $x_1=0$, temos:
        \begin{center}
            \begin{tabular}{lll}
                \toprule
                $x_2$ & $x_3$ & $\log \phi_C(0, x_2, x_3) + \lambda_{x_2 \to \phi_C}(x_2) + \lambda_{x_3 \to \phi_C}(x_3)$\\
                \midrule
                0 & 0 & $\log 4 + \log 8 = \log 32$  \\
                1 & 0 & $\log 2 + \log 8 = \log 16$ \\
                0 & 1 & $\log 2 + \log 4 = \log 8$  \\
                1 & 1 & $\log 6 + \log 4 = \log 24$ \\
                \bottomrule
            \end{tabular}
        \end{center}
        O valor m\'aximo \'e $\log 32$ e, para o backtracking, tamb\'em
        precisamos manter o controle do $\argmax$, que aqui \'e
        $\hat{x}_2=\hat{x}_3=0$.
        
        Para $x_1=1$, temos:
        \begin{center}
            \begin{tabular}{lll}
                \toprule
                $x_2$ & $x_3$ & $\log \phi_C(1, x_2, x_3) + \lambda_{x_2 \to \phi_C}(x_2) + \lambda_{x_3 \to \phi_C}(x_3)$\\
                \midrule
                0 & 0 & $\log 2 + \log 8 = \log 16 $  \\
                1 & 0 & $\log 6 + \log 8 = \log 48$\\
                0 & 1 & $\log 6 + \log 4 = \log 24$ \\
                1 & 1 & $\log 4 + \log 4 = \log 16$ \\
                \bottomrule
            \end{tabular}
        \end{center}
        O valor m\'aximo \'e $\log 48$ e o $\argmax$ \'e $(\hat{x}_2=1, \hat{x}_3=0)$.
        
        Assim, no geral, temos:
        \begin{equation}
            \lambdab_{\phi_C \to x_1} = \begin{pmatrix}
                \lambda_{\phi_C \to x_1}(0)\\
                \lambda_{\phi_C \to x_1}(1)
            \end{pmatrix}
            = \begin{pmatrix}
                \log 32\\
                \log 48
            \end{pmatrix}
        \end{equation}
        e a fun\'c\~ao de back-tracking \'e:
        \begin{equation}
            \lambda^*_{\phi_C \to x_1}(x_1) = \begin{cases}
                (\hat{x}_2 = 0, \hat{x}_3 = 0) & \text{se } x_1=0\\
                (\hat{x}_2=1, \hat{x}_3=0) & \text{se } x_1=1
            \end{cases}
        \end{equation}
        Agora temos todas as mensagens de entrada para o n\'o raiz designado $x_1$. \emph{Ignorando a constante de normaliza\'c\~ao}, obtemos:
        \begin{align}
            \boldsymbol{\gamma}&=\begin{pmatrix}
                \gamma^*(0) \\
                \gamma^*(1)
            \end{pmatrix}
            =  \lambdab_{\phi_A \to x_1} +  \lambdab_{\phi_C \to x_1}\\
            &=  \begin{pmatrix}
                \log 2\\
                \log 4
            \end{pmatrix} +
            \begin{pmatrix}
                \log 32\\
                \log 48
            \end{pmatrix}
            =
            \begin{pmatrix}
                \log 64\\
                \log 192
            \end{pmatrix}
        \end{align}
        O valor $x_1$ para o qual $\gamma^*(x_1)$ \'e maior \'e, portanto,
        $\hat{x}_1 = 1$. Inserindo $\hat{x}_1 = 1$ na fun\'c\~ao de backtracking
        $\lambda^*_{\phi_C \to x_1}(x_1)$, obtemos $(\hat{x}_2=1, \hat{x}_3=0)$. Portanto:
        \begin{equation}
            (\hat{x}_1, \hat{x}_2, \hat{x}_3) = \argmax_{x_1, x_2, x_3} p(x_1, x_2, x_3 | x_4=0, x_5=0) = (1,1,0).
        \end{equation}
        
        Neste exemplo de baixa dimens\~ao, podemos verificar a solu\'c\~ao
        calculando a fmp n\~ao normalizada para todas as combina\'c\~oes de
        $x_1,x_2,x_3$. Isso \'e feito na tabela seguinte, onde come\'camos
        com a tabela para $\phi_C$ e depois multiplicamos os demais
        fatores $\phi_A$, $\tilde{\phi}$ e $\phi_B$. 
        \begin{center}
            \begin{tabular}{lllllll}
                \toprule
                $x_1$ & $x_2$ & $x_3$ & $\phi_C$ & $\phi_C \phi_A$ & $\phi_C \phi_A\tilde{\phi}$& $\phi_C \phi_A\tilde{\phi}\phi_B$ \\
                \midrule
                0 & 0 & 0 & 4 & 8 & 16 & $16\cdot 4$\\
                1 & 0 & 0 & 2 & 8 & 16 & $16 \cdot 4$\\
                0 & 1 & 0 & 2 & 4 & 8 & $8\cdot 4$\\
                1 & 1 & 0 & 6 & 24& 48 & $48 \cdot 4$\\
                0 & 0 & 1 & 2 & 4 & 4  & $4 \cdot 4$\\
                1 & 0 & 1 & 6 & 24 & 24 & $24 \cdot 4$\\
                0 & 1 & 1 & 6 & 12 & 12 & $12 \cdot 4$\\
                1 & 1 & 1 & 4 & 16 & 16 & $16\cdot 4$\\
                \bottomrule
            \end{tabular}
        \end{center}
        Por exemplo, para a coluna $\phi_c \phi_A$, multiplicamos cada
        valor de $\phi_C(x_1, x_2, x_3)$ por $\phi_A(x_1)$, de modo que as
        linhas com $x_1=0$ s\~ao multiplicadas por 2, e as linhas com $x_1=1$
        por 4.
        
        O valor m\'aximo na coluna final \'e alcan\'cado para $x_1=1,
        x_2=1, x_3=0$, em linha com o resultado acima (e $48\cdot 4 =
        192$). Como $\phi_B(x_2)$ \'e uma constante, sendo igual a 4 para todos
        os valores de $x_2$, poder\'iamos t\^e-la ignorado no c\'alculo. A
        raz\~ao formal para isso \'e que, como o modelo n\~ao est\'a normalizado,
        temos permiss\~ao para reescalar cada fator por uma \emph{constante}
        arbitr\'aria (dependente do fator). Esta opera\'c\~ao n\~ao altera o
        modelo. Assim, poder\'iamos dividir $\phi_B$ por 4, o que resultaria em
        um valor de 1, de modo que o fator pode, de fato, ser ignorado.
        
    \end{solution}
    \item Calcule $\argmax_{x_1, \ldots, x_5} p(x_1, \ldots, x_5)$ via passagem de mensagens max-sum com $x_1$ como raiz.
    \label{q:x_1-root-alt}
    \begin{solution}
        Como discutido na solu\'c\~ao acima, podemos descartar o
        fator $\phi_B(x_2)$, j\'a que ele assume o mesmo valor para todos os
        $x_2$. Al\'em disso, podemos reescalar os fatores individuais por uma
        constante para que sejam mais f\'aceis de calcular \`a m\~ao. N\'os os
        normalizamos para que o maior valor seja um, o que resulta nos
        seguintes fatores. Note que isso \'e inteiramente opcional.
        \begin{center}
            \begin{tabular}{ll}
                \toprule
                $x_1$ & $\phi_A$\\
                \midrule
                0 & 1\\
                1 & 2\\
                \bottomrule
            \end{tabular}
            \hfill
            \begin{tabular}{llll}
                \toprule
                $x_1$ & $x_2$ & $x_3$ & $\phi_C$\\
                \midrule
                0 & 0 & 0 & 2 \\
                1 & 0 & 0 & 1 \\
                0 & 1 & 0 & 1 \\
                1 & 1 & 0 & 3 \\
                0 & 0 & 1 & 1 \\
                1 & 0 & 1 & 3 \\
                0 & 1 & 1 & 3 \\
                1 & 1 & 1 & 2 \\
                \bottomrule
            \end{tabular}
            \hfill
            \begin{tabular}{lll}
                \toprule
                $x_3$ & $x_4$ & $\phi_D$\\
                \midrule
                0 & 0 &  4 \\
                1 & 0 &  1 \\
                0 & 1 &  1 \\
                1 & 1 &  3 \\
                \bottomrule
            \end{tabular}
            \hfill
            \begin{tabular}{lll}
                \toprule
                $x_3$ & $x_5$ & $\phi_E$\\
                \midrule
                0 & 0 &  1 \\
                1 & 0 &  2 \\
                0 & 1 &  2 \\
                1 & 1 &  1 \\
                \bottomrule
            \end{tabular}
            \hfill
            \begin{tabular}{ll}
                \toprule
                $x_5$ & $\phi_F$\\
                \midrule
                0 & 1\\
                1 & 8\\
                \bottomrule
            \end{tabular}   
        \end{center}
        O grafo de fatores sem $\phi_B$, juntamente com as mensagens que precisamos calcular, \'e:
        \begin{center}
            \begin{tikzpicture}[ugraph]
                \node[fact, label=above: $\phi_A$] (fa) at (0,0) {} ; %
                \node[cont] (x1) at (1.5,0)  {$x_1$} ; %
                \node[fact, label=below: $\phi_C$] (fc) at (3,0) {} ; %
                \node[cont] (x2) at (3,1)  {$x_2$} ; %
                \node[cont] (x3) at (4.5,0)  {$x_3$} ; %
                \node[fact, label=above: $\phi_D$] (fd) at (5.5,1) {} ; %
                \node[cont] (x4) at (7,1)  {$x_4$} ; %
                \node[fact, label=above: $\phi_E$] (fe) at (5.5,-1) {} ; %
                \node[cont] (x5) at (7,-1)  {$x_5$} ; %
                \node[fact, label=above: $\phi_F$] (ff) at (8.5,-1) {} ; %
                \draw (fa) -- (x1) node[midway,above] {$\rightarrow$}; 
                \draw (x1) -- (fc) node[midway,above] {$\leftarrow$}; 
                \draw (fc) -- (x2) node[midway,left] {$\downarrow$}; 
                \draw (fc) -- (x3) node[midway,above] {$\leftarrow$}; 
                \draw (x3) -- (fd) node[midway,above, sloped] {$\leftarrow$}; 
                \draw (x3) -- (fe) node[midway,below, sloped] {$\leftarrow$}; 
                \draw (fd) -- (x4) node[midway,above, sloped] {$\leftarrow$}; 
                \draw (fe) -- (x5) node[midway,below, sloped] {$\leftarrow$}; 
                \draw (x5) -- (ff) node[midway,below, sloped] {$\leftarrow$}; 
            \end{tikzpicture}
        \end{center}
        As mensagens das folhas (log) s\~ao (usando a nota\'c\~ao vetorial onde o elemento superior corresponde a $x_i=0$ e o inferior a $x_i=1$):
        \begin{align}
            \lambdab_{\phi_A \to x_1} & = \begin{pmatrix}
                0\\
                \log 2
            \end{pmatrix}
            &
            \lambdab_{x_2 \to \phi_C} & = \begin{pmatrix}
                0\\
                0
            \end{pmatrix}
            &
            \lambdab_{x_4 \to \phi_D} &  = \begin{pmatrix}
                0\\
                0
            \end{pmatrix}
            &
            \lambdab_{\phi_F \to x_5} &  = \begin{pmatrix}
                0\\
                \log 8
            \end{pmatrix}
        \end{align}
        O n\'o de vari\'avel $x_5$ possui apenas uma aresta de entrada, de modo que
        $\lambdab_{x_5 \to \phi_E} = \lambdab_{\phi_F \to x_5}$. A
        mensagem $\lambda_{\phi_E \to x_3}(x_3)$ \'e igual a:
        \begin{align}
            \lambda_{\phi_E \to x_3}(x_3) & = \max_{x_5} \{ \log \phi_E(x_3, x_5) + \lambda_{x_5 \to \phi_E}(x_5) \}
        \end{align}
        Escrevendo $\log \phi_E(x_3, x_5) + \lambda_{x_5 \to \phi_E}(x_5)$ para todos os $x_5$ em fun\'c\~ao de $x_3$, temos:
        \begin{center}
            \begin{tabular}{ll}
                \toprule
                $x_5$ &  $\log \phi_E(0, x_5) + \lambda_{x_5 \to \phi_E}(x_5)$\\
                \midrule
                0  &  $\log 1 + 0 = 0$\\
                1  &  $\log 2 + \log 8 = \log 16$\\
                \bottomrule
            \end{tabular}
            \hspace{3ex}
            \begin{tabular}{ll}
                \toprule
                $x_5$ &  $\log \phi_E(1, x_5) + \lambda_{x_5 \to \phi_E}(x_5)$\\
                \midrule
                0  &  $\log 2 + 0 = \log 2$\\
                1  &  $\log 1 + \log 8 = \log 8$\\
                \bottomrule
            \end{tabular}
        \end{center}
        Tomando o m\'aximo sobre $x_5$ como uma fun\'c\~ao de $x_3$, obtemos:
        \begin{equation}
            \lambdab_{\phi_E \to x_3} = \begin{pmatrix}
                \log 16\\
                \log 8
            \end{pmatrix}
        \end{equation}
        e a fun\'c\~ao de backtracking que indica o maximizador
        $\hat{x}_5 = \argmax_{x_5} \{ \log \phi_E(x_3, x_5) + \lambda_{x_5 \to \phi_E}(x_5) \}$ em fun\'c\~ao de $x_3$ \'e igual a:
        \begin{equation}
            \lambda^*_{\phi_E \to x_3}(x_3) = \begin{cases}
                \hat{x}_5=1 & \text{se } x_3= 0\\
                \hat{x}_5=1 & \text{se } x_3= 1
            \end{cases}
        \end{equation}
        Realizamos o mesmo tipo de opera\'c\~ao para $\lambda_{\phi_D \to x_3}(x_3)$:
        \begin{align}
            \lambda_{\phi_D \to x_3}(x_3) & = \max_{x_4} \{ \log \phi_D(x_3, x_4) + \lambda_{x_4 \to \phi_D}(x_4) \}
        \end{align}
        Como $\lambda_{x_4 \to \phi_D}(x_4)=0$ para todo $x_4$, a tabela com todos os valores de $\log \phi_D(x_3, x_4) + \lambda_{x_4 \to \phi_D}(x_4)$ \'e:
        \begin{center}
            \begin{tabular}{lll}
                \toprule
                $x_3$ & $x_4$ & $\log \phi_D(x_3, x_4) + \lambda_{x_4 \to \phi_D}(x_4)$ \\
                \midrule
                0 & 0 &  $\log 4 + 0 = \log 4$ \\
                1 & 0 &  $\log 1 + 0 = 0$\\
                0 & 1 &  $\log 1 + 0 = 0$\\
                1 & 1 &  $\log 3 + 0 = \log 3$ \\
                \bottomrule
            \end{tabular}
        \end{center}
        Tomando o m\'aximo sobre $x_4$ como uma fun\'c\~ao de $x_3$, obtemos portanto:
        \begin{equation}
            \lambdab_{\phi_D \to x_3} = \begin{pmatrix}
                \log 4\\
                \log 3
            \end{pmatrix}
        \end{equation}
        e a fun\'c\~ao de backtracking que indica o maximizador
        $\hat{x}_4 = \argmax_{x_4} \{ \log \phi_D(x_3, x_4) + \lambda_{x_4 \to \phi_D}(x_4) \}$ em fun\'c\~ao de $x_3$ \'e igual a:
        \begin{equation}
            \lambda^*_{\phi_D \to x_3}(x_3) = \begin{cases}
                \hat{x}_4=0 & \text{se } x_3= 0\\
                \hat{x}_4=1 & \text{se } x_3= 1
            \end{cases}
        \end{equation}
        Para a mensagem $\lambda_{x_3 \to \phi_C}(x_3)$, somamos as
        mensagens $\lambda_{\phi_E \to x_3}(x_3)$ e
        $\lambda_{\phi_D \to x_3}(x_3)$, o que resulta em:
        \begin{equation}
            \lambdab_{ x_3 \to \phi_C} = \begin{pmatrix}
                \log 16 + \log 4\\
                \log 8 + \log 3\\
            \end{pmatrix}
            =\begin{pmatrix}
                \log 64\\
                \log 24\\
            \end{pmatrix}
        \end{equation}
        
        Em seguida, calculamos a mensagem $\lambda_{\phi_C \to x_1}(x_1)$ ao maximizar sobre $x_2$ e $x_3$:
        \begin{equation}
            \lambda_{\phi_C \to x_1}(x_1)  = \max_{x_2, x_3} \{ \log \phi_C(x_1, x_2, x_3) + \lambda_{x_2 \to \phi_C}(x_2) + \lambda_{x_3 \to \phi_C}(x_3) \}
        \end{equation}
        Como $\lambda_{x_2 \to \phi_C}(x_2) = 0$, o problema torna-se:
        \begin{equation}
            \lambda_{\phi_C \to x_1}(x_1)  = \max_{x_2, x_3} \{ \log \phi_C(x_1, x_2, x_3) + \lambda_{x_3 \to \phi_C}(x_3) \}
        \end{equation}
        Baseando-nos na tabela para $\phi_C$, formamos uma tabela com todos
        os valores de $\log \phi_C(x_1, x_2, x_3) + \lambda_{x_3 \to \phi_C}(x_3)$:
        \begin{center}
            \begin{tabular}{llll}
                \toprule
                $x_1$ & $x_2$ & $x_3$ & $\log \phi_C(x_1, x_2, x_3) + \lambda_{x_3 \to \phi_C}(x_3)$\\
                \midrule
                0 & 0 & 0 & $\log 2 + \log 64 = \boldsymbol{\log 128}$ \\
                1 & 0 & 0 & $0 + \log 64 = \log 64$ \\
                0 & 1 & 0 & $0 + \log 64 = \log 64$ \\
                1 & 1 & 0 & $\log 3 + \log 64 = \boldsymbol{\log 192}$ \\
                0 & 0 & 1 & $0 + \log 24 = \log 24$ \\
                1 & 0 & 1 & $\log 3 + \log 24 = \log 72$ \\
                0 & 1 & 1 & $\log 3 + \log 24 = \log 72 $ \\
                1 & 1 & 1 & $\log 2 + \log 24 = \log 48$ \\
                \bottomrule
            \end{tabular}
        \end{center}
        Os valores m\'aximos como uma fun\'c\~ao de $x_1$ est\~ao destacados na tabela, o que fornece a mensagem:
        \begin{equation}
            \lambdab_{\phi_C \to x_1} = \begin{pmatrix}
                \log 128\\
                \log 192
            \end{pmatrix}
        \end{equation}
        e a fun\'c\~ao de backtracking:
        \begin{equation}
            \lambda^*_{\phi_C \to x_1}(x_1) = \begin{cases}
                (\hat{x}_2=0, \hat{x}_3=0) & \text{se } x_1=0\\
                (\hat{x}_2=1, \hat{x}_3=0) & \text{se } x_1=1
            \end{cases}
        \end{equation}
        Agora temos todas as mensagens de entrada para o n\'o raiz designado $x_1$. \emph{Ignorando a constante de normaliza\'c\~ao}, obtemos:
        \begin{align}
            \boldsymbol{\gamma}&=\begin{pmatrix}
                \gamma^*(0) \\
                \gamma^*(1)
            \end{pmatrix}
            = \begin{pmatrix}
                0 + \log 128 \\
                \log 2 + \log 192 
            \end{pmatrix}
        \end{align}
        Podemos agora iniciar o backtracking para calcular o desejado
        $\argmax_{x_1, \ldots, x_5} p(x_1, \ldots, x_5)$. Come\'cando na
        raiz, temos $\hat{x}_1 = \argmax_{x_1} \gamma^*(x_1) = 1$. Inserindo este valor na tabela de consulta $\lambda^*_{\phi_C \to x_1}(x_1)$, obtemos $(\hat{x}_2=1, \hat{x}_3=0)$. Com a tabela de consulta $\lambda^*_{\phi_E \to x_3}(0)$, encontramos $\hat{x}_5=1$ e $\lambda^*_{\phi_D \to x_3}(0)$ fornece $\hat{x}_4=0$, de modo que, no geral:
        \begin{equation}
            \argmax_{x_1, \ldots, x_5} p(x_1, \ldots, x_5) = (1, 1, 0, 0, 1).
        \end{equation}
        
        
    \end{solution}
    
    \item Calcule $\argmax_{x_1, \ldots, x_5} p(x_1, \ldots, x_5)$ via passagem de mensagens max-sum com $x_3$ como raiz.
    \begin{solution}
        Com $x_3$ como raiz, precisamos das seguintes mensagens:
        \begin{center}
            \begin{tikzpicture}[ugraph]
                \node[fact, label=above: $\phi_A$] (fa) at (0,0) {} ; %
                \node[cont] (x1) at (1.5,0)  {$x_1$} ; %
                \node[fact, label=below: $\phi_C$] (fc) at (3,0) {} ; %
                \node[cont] (x2) at (3,1)  {$x_2$} ; %
                \node[cont] (x3) at (4.5,0)  {$x_3$} ; %
                \node[fact, label=above: $\phi_D$] (fd) at (5.5,1) {} ; %
                \node[cont] (x4) at (7,1)  {$x_4$} ; %
                \node[fact, label=above: $\phi_E$] (fe) at (5.5,-1) {} ; %
                \node[cont] (x5) at (7,-1)  {$x_5$} ; %
                \node[fact, label=above: $\phi_F$] (ff) at (8.5,-1) {} ; %
                \draw (fa) -- (x1) node[midway,above] {$\rightarrow$}; 
                \draw (x1) -- (fc) node[midway,above] {$\rightarrow$}; 
                \draw (fc) -- (x2) node[midway,left] {$\downarrow$}; 
                \draw (fc) -- (x3) node[midway,above] {$\rightarrow$}; 
                \draw (x3) -- (fd) node[midway,above, sloped] {$\leftarrow$}; 
                \draw (x3) -- (fe) node[midway,below, sloped] {$\leftarrow$}; 
                \draw (fd) -- (x4) node[midway,above, sloped] {$\leftarrow$}; 
                \draw (fe) -- (x5) node[midway,below, sloped] {$\leftarrow$}; 
                \draw (x5) -- (ff) node[midway,below, sloped] {$\leftarrow$}; 
            \end{tikzpicture}
        \end{center}
        As seguintes mensagens s\~ao as mesmas de quando $x_1$ era a raiz:
        \begin{align}
            \lambdab_{\phi_D \to x_3} &= \begin{pmatrix}
                \log 4\\
                \log 3
            \end{pmatrix}
            &
            \lambdab_{\phi_E \to x_3} &= \begin{pmatrix}
                \log 16\\
                \log 8
            \end{pmatrix}
            &
            \lambdab_{\phi_A \to x_1} & = \begin{pmatrix}
                0\\
                \log 2
            \end{pmatrix}
            &
            \lambdab_{x_2 \to \phi_C} & = \begin{pmatrix}
                0\\
                0
            \end{pmatrix}
        \end{align}
        Como $x_1$ possui apenas uma mensagem de entrada, temos ainda:
        \begin{equation}
            \lambdab_{x_1 \to \phi_C} = \lambdab_{\phi_A \to x_1} =  \begin{pmatrix}
                0\\
                \log 2
            \end{pmatrix}.
        \end{equation}
        Em seguida, computamos $\lambda_{\phi_C \to x_3}(x_3)$,
        \begin{equation}
            \lambda_{\phi_C \to x_3}(x_3) = \max_{x_1, x_2} \{ \log \phi_C(x_1,x_2,x_3) + \lambda_{x_1 \to \phi_C}(x_1) + \lambda_{x_2 \to \phi_C}(x_2) \}.
        \end{equation}
        Primeiro formamos uma tabela para $\log \phi_C(x_1,x_2,x_3) + \lambda_{x_1 \to \phi_C}(x_1) + \lambda_{x_2 \to \phi_C}(x_2)$, notando que $\lambda_{x_2 \to \phi_C}(x_2)=0$:
        \begin{center}
            \begin{tabular}{llll}
                \toprule
                $x_1$ & $x_2$ & $x_3$ & $\log \phi_C(x_1,x_2,x_3) + \lambda_{x_1 \to \phi_C}(x_1) + \lambda_{x_2 \to \phi_C}(x_2)$  \\
                \midrule
                0 & 0 & 0 & $\log 2 + 0 = \log 2$ \\
                1 & 0 & 0 & $0+\log 2 = \log 2$ \\
                0 & 1 & 0 & $0+0 = 0$ \\
                1 & 1 & 0 & $\log 3+ \log 2 = \boldsymbol{\log 6}$ \\
                0 & 0 & 1 & $0+0 = 0 $ \\
                1 & 0 & 1 & $\log 3+ \log 2 = \boldsymbol{\log 6}$ \\
                0 & 1 & 1 & $\log 3+0 = \log 3$ \\
                1 & 1 & 1 & $\log 2+\log 2 = \log 4$ \\
                \bottomrule
            \end{tabular}
        \end{center}
        Os valores m\'aximos em fun\'c\~ao de $x_3$ est\~ao destacados na tabela, o que fornece a mensagem:
        \begin{equation}
            \lambdab_{\phi_C \to x_3} = \begin{pmatrix}
                \log 6\\
                \log 6
            \end{pmatrix}
        \end{equation}
        e a fun\'c\~ao de backtracking:
        \begin{equation}
            \lambda^*_{\phi_C \to x_3}(x_3) = \begin{cases}
                (\hat{x}_1=1, \hat{x}_2=1) & \text{se } x_3=0\\
                (\hat{x}_1=1, \hat{x}_2=0) & \text{se } x_3=1
            \end{cases}
        \end{equation}
        Agora temos todas as mensagens de entrada para $x_3$ e podemos computar
        $\gamma^*(x_3)$ at\'e a constante de normaliza\'c\~ao $-\log Z$ (que n\~ao \'e
        necess\'aria se estivermos interessados apenas no $\argmax$):
        \begin{align}
            \boldsymbol{\gamma}&=\begin{pmatrix}
                \gamma^*(0) \\
                \gamma^*(1)
            \end{pmatrix}
            =  \lambdab_{\phi_C \to x_3} +  \lambdab_{\phi_D \to x_3} +  \lambdab_{\phi_E \to x_3}\\
            & = \begin{pmatrix}
                \log 6 + \log 4 + \log 16 = \log 384 \\
                \log 6 + \log 3 + \log 8 = \log 144
            \end{pmatrix}
        \end{align}
        Podemos agora iniciar o backtracking, que resulta em: $\hat{x}_3 = 0$,
        de modo que $\lambda^*_{\phi_C \to x_3}(0) = (\hat{x}_1=1, \hat{x}_2=1)$. As fun\'c\~oes de backtracking $\lambda^*_{\phi_E \to x_3}(x_3)$ e $\lambda^*_{\phi_D \to x_3}(x_3)$ s\~ao as mesmas da quest\~ao \ref{q:x_1-root-alt}, o que resulta em $\lambda^*_{\phi_E \to x_3}(0) = \hat{x}_5=1$ e $\lambda^*_{\phi_D \to x_3}(0) = \hat{x}_4=0$. Assim, no geral, descobrimos que:
        \begin{equation}
            \argmax_{x_1, \ldots, x_5} p(x_1, \ldots, x_5) = (1, 1, 0, 0, 1).
        \end{equation}
        Note que isto corresponde ao resultado da quest\~ao \ref{q:x_1-root-alt} onde $x_1$ era a raiz. Isso ocorre porque a sa\'ida do algoritmo max-sum \'e invariante \`a escolha da raiz.
        
    \end{solution}
    
\end{exenumerate}

\ex{Escolha da ordem de elimina\'c\~ao em grafos de fatores}
Considere o seguinte grafo de fatores, que cont\'em um loop:

\begin{center}
    \begin{tikzpicture}[ugraph]
        \node[cont] (x1) at (0,0)  {$x_1$} ; %
        \node[fact, label=above: $\phi_A$] (fa) at (1.5,0) {} ; %
        \node[cont] (x2) at (2.5,1)  {$x_2$} ; %
        \node[cont] (x3) at (2.5,-1)  {$x_3$} ; %
        \node[fact, label=above: $\phi_B$] (fb) at (3.5,0) {} ; %
        \node[cont] (x4) at (5,0)  {$x_4$} ; %
        \node[fact, label=above: $\phi_C$] (fc) at (6,1) {} ; %
        \node[cont] (x5) at (7.5,1)  {$x_5$} ; %
        \node[cont] (x6) at (7.5,-1)  {$x_6$} ;
        \node[fact, label=above: $\phi_D$] (fd) at (6,-1) {} ; %
        
        \draw (fd) -- (x6);
        \draw (x1) -- (fa);
        \draw (fa) -- (x2);
        \draw (fa) -- (x3);
        \draw (x2) -- (fb);
        \draw (x3) -- (fb);
        \draw (fb) -- (x4);
        \draw (x4) -- (fc);
        \draw (x4) -- (fd);
        \draw (fc) -- (x5);
    \end{tikzpicture}
\end{center}

Sejam todas as vari\'aveis bin\'arias, $x_i \in \{0,1\}$, e os fatores definidos como segue:
\begin{center}
    \hfill
    %
    % phi_A (x_1, x_2, x_3)
    \begin{tabular}{llll}
        \toprule
        $x_1$ & $x_2$ & $x_3$ & $\phi_A$\\
        \midrule
        0 & 0 & 0 & 4 \\
        1 & 0 & 0 & 2 \\
        0 & 1 & 0 & 2 \\
        1 & 1 & 0 & 6 \\
        0 & 0 & 1 & 2 \\
        1 & 0 & 1 & 6 \\
        0 & 1 & 1 & 6 \\
        1 & 1 & 1 & 4 \\
        \bottomrule
    \end{tabular}
    \hfill
    %
    % phi_B (x_2, x_3, x_4)
    \begin{tabular}{llll}
        \toprule
        $x_2$ & $x_3$ & $x_4$ & $\phi_B$\\
        \midrule
        0 & 0 & 0 & 2 \\
        1 & 0 & 0 & 2 \\
        0 & 1 & 0 & 4 \\
        1 & 1 & 0 & 2 \\
        0 & 0 & 1 & 6 \\
        1 & 0 & 1 & 8 \\
        0 & 1 & 1 & 4 \\
        1 & 1 & 1 & 2 \\
        \bottomrule
    \end{tabular}
    \hfill
    %
    % phi_C (x_4, x_5)
    \begin{tabular}{lll}
        \toprule
        $x_4$ & $x_5$ & $\phi_C$\\
        \midrule
        0 & 0 & 8 \\
        1 & 0 & 2 \\
        0 & 1 & 2 \\
        1 & 1 & 6 \\
        \bottomrule
    \end{tabular}
    \hfill
    %
    % phi_D (x_4, x_6)
    \begin{tabular}{lll}
        \toprule
        $x_4$ & $x_6$ & $\phi_D$\\
        \midrule
        0 & 0 & 3 \\
        1 & 0 & 6 \\
        0 & 1 & 6 \\
        1 & 1 & 3 \\
        \bottomrule
    \end{tabular}
    \hfill
\end{center}

\begin{exenumerate}
    
    %
    % Question a – condition on x_1 and x_6.
    \item Desenhe o grafo de fatores correspondente a $p(x_2, x_3, x_4, x_5 \mid x_1=0, x_6=1)$ e forne\'ca as tabelas que definem os novos fatores $\phi_A^{x_1=0}(x_2, x_3)$ e $\phi_D^{x_6=1}(x_4)$ que voc\^e obteve.
    
    %
    % Question a – solution
    \begin{solution}
        Primeiro condicione em $x_1 = 0$:
        
        O n\'o de fator $\phi_A(x_1, x_2, x_3)$ depende de $x_1$, portanto criamos um novo fator $\phi_A^{x_1=0}(x_2, x_3)$ a partir da tabela de $\phi_A$ usando as linhas onde $x_1 = 0$.
        
        % p(x_2, x_3, x_4, x_5, x_6 | x_1=0)
        \begin{center}
            \begin{tikzpicture}[ugraph]
                \node[fact, label=above: $\phi_A^{x_1=0}$] (fa) at (1.5,0) {} ; %
                \node[cont] (x2) at (2.5,1)  {$x_2$} ; %
                \node[cont] (x3) at (2.5,-1)  {$x_3$} ; %
                \node[fact, label=above: $\phi_B$] (fb) at (3.5,0) {} ; %
                \node[cont] (x4) at (5,0)  {$x_4$} ; %
                \node[fact, label=above: $\phi_C$] (fc) at (6,1) {} ; %
                \node[cont] (x5) at (7.5,1)  {$x_5$} ; %
                \node[cont] (x6) at (7.5,-1)  {$x_6$} ;
                \node[fact, label=above: $\phi_D$] (fd) at (6,-1) {} ; %
                
                \draw (fd) -- (x6);
                \draw (fa) -- (x2);
                \draw (fa) -- (x3);
                \draw (x2) -- (fb);
                \draw (x3) -- (fb);
                \draw (fb) -- (x4);
                \draw (x4) -- (fc);
                \draw (x4) -- (fd);
                \draw (fc) -- (x5);
            \end{tikzpicture}
        \end{center}
        
        % phi_A^{x_1=0} (x_2, x_3)
        \begin{center}
            %
            % phi_A (x_1, x_2, x_3)
            \begin{tabular}{lllll}
                \toprule
                & $x_1$ & $x_2$ & $x_3$ & $\phi_A$\\
                \midrule
                $\rightarrow$ & 0 & 0 & 0 & 4 \\
                & 1 & 0 & 0 & 2 \\
                $\rightarrow$ & 0 & 1 & 0 & 2 \\
                & 1 & 1 & 0 & 6 \\
                $\rightarrow$ & 0 & 0 & 1 & 2 \\
                & 1 & 0 & 1 & 6 \\
                $\rightarrow$ & 0 & 1 & 1 & 6 \\
                & 1 & 1 & 1 & 4 \\
                \bottomrule
            \end{tabular}
            \hspace{3ex} \text{de modo que} \hspace{3ex}
            %
            % phi_A (x_2, x_3) for x_1 = 0
            \begin{tabular}{lll}
                \toprule
                $x_2$ & $x_3$ & $\phi_A^{x_1=0}$\\
                \midrule
                0 & 0 & 4 \\
                1 & 0 & 2 \\
                0 & 1 & 2 \\
                1 & 1 & 6 \\
                \bottomrule
            \end{tabular}
        \end{center}
        
        Em seguida, condicione em $x_6 = 1$:
        
        O n\'o de fator $\phi_D(x_4, x_6)$ depende de $x_6$, portanto criamos um novo fator $\phi_D^{x_6=1}(x_4)$ a partir da tabela de $\phi_D$ usando as linhas onde $x_6 = 1$.
        
        % p(x_2, x_3, x_4, x_5 | x_1=0, x_6=1)$
        \begin{center}
            \begin{tikzpicture}[ugraph]
                \node[fact, label=above: $\phi_A^{x_1=0}$] (fa) at (1.5,0) {} ; %
                \node[cont] (x2) at (2.5,1)  {$x_2$} ; %
                \node[cont] (x3) at (2.5,-1)  {$x_3$} ; %
                \node[fact, label=above: $\phi_B$] (fb) at (3.5,0) {} ; %
                \node[cont] (x4) at (5,0)  {$x_4$} ; %
                \node[fact, label=above: $\phi_C$] (fc) at (6,1) {} ; %
                \node[cont] (x5) at (7.5,1)  {$x_5$} ; %
                \node[fact, label=right: $\phi_D^{x_6=1}$] (fd) at (6,-1) {} ; %
                
                \draw (fa) -- (x2);
                \draw (fa) -- (x3);
                \draw (x2) -- (fb);
                \draw (x3) -- (fb);
                \draw (fb) -- (x4);
                \draw (x4) -- (fc);
                \draw (x4) -- (fd);
                \draw (fc) -- (x5);
            \end{tikzpicture}
        \end{center}
        
        % phi_D^{x_6=1} (x_4)
        \begin{center}
            %
            % phi_D (x_4, x_6)
            \begin{tabular}{llll}
                \toprule
                & $x_4$ & $x_6$ & $\phi_D$\\
                \midrule
                & 0 & 0 & 3 \\
                & 1 & 0 & 6 \\
                $\rightarrow$ & 0 & 1 & 6 \\
                $\rightarrow$ & 1 & 1 & 3 \\
                \bottomrule
            \end{tabular}
            \hspace{3ex} \text{de modo que} \hspace{3ex}
            %
            % phi_D (x_4) for x_6 = 1
            \begin{tabular}{ll}
                \toprule
                $x_4$ & $\phi_D^{x_6=1}$\\
                \midrule
                0 & 6 \\
                1 & 3 \\
                \bottomrule
            \end{tabular}
        \end{center}
        
    \end{solution}
    
    %
    % Question b – find p(x_2 | x_1, x_6) with a bad variable ordering.
    \item Encontre $p(x_2 \mid x_1=0, x_6=1)$ usando a ordem de elimina\'c\~ao $(x_4, x_5, x_3)$:
    
    \begin{exenumerate}
        \item Desenhe o grafo para $p(x_2, x_3, x_5 \mid x_1=0, x_6=1)$ ao marginalizar $x_4$ \\
        Calcule a tabela para o novo fator $\tilde{\phi}_4(x_2, x_3, x_5)$ \\
        
        \item Desenhe o grafo para $p(x_2, x_3 \mid x_1=0, x_6=1)$ ao marginalizar $x_5$ \\
        Calcule a tabela para o novo fator $\tilde{\phi}_{45}(x_2, x_3)$ \\
        
        \item Desenhe o grafo para $p(x_2 \mid x_1=0, x_6=1)$ ao marginalizar $x_3$ \\
        Calcule a tabela para o novo fator $\tilde{\phi}_{453}(x_2)$ 
    \end{exenumerate}
    
    %
    % Question b – solution
    \begin{solution}
        
        Come\'cando com o grafo de fatores para $p(x_2, x_3, x_4, x_5 \mid x_1=0, x_6=1)$:
        
        \begin{center}
            \begin{tikzpicture}[ugraph]
                \node[fact, label=above: $\phi_A^{x_1=0}$] (fa) at (1.5,0) {} ; %
                \node[cont] (x2) at (2.5,1)  {$x_2$} ; %
                \node[cont] (x3) at (2.5,-1)  {$x_3$} ; %
                \node[fact, label=above: $\phi_B$] (fb) at (3.5,0) {} ; %
                \node[cont] (x4) at (5,0)  {$x_4$} ; %
                \node[fact, label=above: $\phi_C$] (fc) at (6,1) {} ; %
                \node[cont] (x5) at (7.5,1)  {$x_5$} ; %
                \node[fact, label=right: $\phi_D^{x_6=1}$] (fd) at (6,-1) {} ; %
                
                \draw (fa) -- (x2);
                \draw (fa) -- (x3);
                \draw (x2) -- (fb);
                \draw (x3) -- (fb);
                \draw (fb) -- (x4);
                \draw (x4) -- (fc);
                \draw (x4) -- (fd);
                \draw (fc) -- (x5);
            \end{tikzpicture}
        \end{center}
        
        A marginaliza\'c\~ao de $x_4$ combina os tr\^es fatores $\phi_B$, $\phi_C$ e $\phi_D^{x_6=1}$:
        
        \begin{center}
            \begin{tikzpicture}[ugraph]
                \node[fact, label=above: $\phi_A^{x_1=0}$] (fa) at (1.5,0) {} ; %
                \node[cont] (x2) at (2.5,1)  {$x_2$} ; %
                \node[cont] (x3) at (2.5,-1)  {$x_3$} ; %
                \node[fact, label=above: $\tilde{\phi}_4$] (f4) at (3.5,0) {} ; %
                \node[cont] (x5) at (5,0)  {$x_5$} ; %
                
                \draw (fa) -- (x2);
                \draw (fa) -- (x3);
                \draw (x2) -- (f4);
                \draw (x3) -- (f4);
                \draw (f4) -- (x5);
            \end{tikzpicture}
        \end{center}
        
        A marginaliza\'c\~ao de $x_5$ modifica o fator $\tilde{\phi}_4$:
        
        \begin{center}
            \begin{tikzpicture}[ugraph]
                \node[fact, label=above: $\phi_A^{x_1=0}$] (fa) at (1.5,0) {} ; %
                \node[cont] (x2) at (2.5,1)  {$x_2$} ; %
                \node[cont] (x3) at (2.5,-1)  {$x_3$} ; %
                \node[fact, label=right: $\tilde{\phi}_{45}$] (f45) at (3.5,0) {} ; %
                
                \draw (fa) -- (x2);
                \draw (fa) -- (x3);
                \draw (x2) -- (f45);
                \draw (x3) -- (f45);
            \end{tikzpicture}
        \end{center}
        
        A marginaliza\'c\~ao de $x_3$ combina os fatores $\phi_A^{x_1=0}$ e $\tilde{\phi}_{45}$:
        
        \begin{center}
            \begin{tikzpicture}[ugraph]
                \node[cont] (x2) at (0,0)  {$x_2$} ; %
                \node[fact, label=right: $\tilde{\phi}_{453}$] (f453) at (2,0) {} ; %
                
                \draw (x2) -- (f453);
            \end{tikzpicture}
        \end{center}
        
        Agora calculamos as tabelas para os novos fatores $\tilde{\phi}_4$, $\tilde{\phi}_{45}$, $\tilde{\phi}_{453}$.
        
        Primeiro encontre $\tilde{\phi}_4(x_2, x_3, x_5)$:
        \begin{center}
            %
            % phi_B (x_2, x_3, x_4)
            \begin{tabular}{llll}
                \toprule
                $x_2$ & $x_3$ & $x_4$ & $\phi_B$\\
                \midrule
                0 & 0 & 0 & 2 \\
                1 & 0 & 0 & 2 \\
                0 & 1 & 0 & 4 \\
                1 & 1 & 0 & 2 \\
                0 & 0 & 1 & 6 \\
                1 & 0 & 1 & 8 \\
                0 & 1 & 1 & 4 \\
                1 & 1 & 1 & 2 \\
                \bottomrule
            \end{tabular}\hspace{2ex}
            %
            % phi_C (x_4, x_5)
            \begin{tabular}{lll}
                \toprule
                $x_4$ & $x_5$ & $\phi_C$\\
                \midrule
                0 & 0 & 8 \\
                1 & 0 & 2 \\
                0 & 1 & 2 \\
                1 & 1 & 6 \\
                \bottomrule
            \end{tabular}\hspace{2ex}
            %
            % phi_D (x_4) for x_6 = 1
            \begin{tabular}{ll}
                \toprule
                $x_4$ & $\phi_D^{x_6=1}$\\
                \midrule
                0 & 6 \\
                1 & 3 \\
                \bottomrule
            \end{tabular}
        \end{center}
        de modo que $\phi_*(x_2,x_3,x_4,x_5) = \phi_B(x_2, x_3, x_4) \phi_C(x_4, x_5) \phi_D^{x_6=1}(x_4)$ \'e igual a:
        \begin{center}
            \begin{tabular}{lllll}
                \toprule
                $x_2$ & $x_3$ & $x_4$& $x_5$ & $\phi_*(x_2,x_3,x_4,x_5)$\\
                \midrule
                0 & 0 & 0 & 0 & 2 * 8 * 6 \\
                1 & 0 & 0 & 0 & 2 * 8 * 6 \\
                0 & 1 & 0 & 0 & 4 * 8 * 6 \\
                1 & 1 & 0 & 0 & 2 * 8 * 6 \\
                0 & 0 & 1 & 0 & 6 * 2 * 3\\
                1 & 0 & 1 & 0 & 8 * 2 * 3\\
                0 & 1 & 1 & 0 & 4 * 2 * 3\\
                1 & 1 & 1 & 0 & 2 * 2 * 3\\
                0 & 0 & 0 & 1 & 2 * 2 * 6\\
                1 & 0 & 0 & 1 & 2 * 2 * 6\\
                0 & 1 & 0 & 1 & 4 * 2 * 6\\
                1 & 1 & 0 & 1 & 2 * 2 * 6\\
                0 & 0 & 1 & 1 & 6 * 6 * 3\\
                1 & 0 & 1 & 1 & 8 * 6 * 3\\
                0 & 1 & 1 & 1 & 4 * 6 * 3\\
                1 & 1 & 1 & 1 & 2 * 6 * 3\\
                \bottomrule
            \end{tabular}
        \end{center}
        e
        \begin{center}
            \begin{tabular}{llllll}
                \toprule
                $x_2$ & $x_3$ & $x_5$ & $\sum_{x_4} \phi_B(x_2, x_3, x_4) \phi_C(x_4, x_5) \phi_D^{x_6=1}(x_4)$ &  & $\tilde{\phi}_4$\\
                \midrule
                0 & 0 & 0 & (2 * 8 * 6) + (6 * 2 * 3) & = & 132 \\
                1 & 0 & 0 & (2 * 8 * 6) + (8 * 2 * 3) & = & 144 \\
                0 & 1 & 0 & (4 * 8 * 6) + (4 * 2 * 3) & = & 216 \\
                1 & 1 & 0 & (2 * 8 * 6) + (2 * 2 * 3) & = & 108 \\
                0 & 0 & 1 & (2 * 2 * 6) + (6 * 6 * 3) & = & 132 \\
                1 & 0 & 1 & (2 * 2 * 6) + (8 * 6 * 3) & = & 168 \\
                0 & 1 & 1 & (4 * 2 * 6) + (4 * 6 * 3) & = & 120 \\
                1 & 1 & 1 & (2 * 2 * 6) + (2 * 6 * 3) & = & 60 \\
                \bottomrule
            \end{tabular}
        \end{center}
        
        Em seguida, encontre $\tilde{\phi}_{45}(x_2, x_3)$:
        \begin{center}
            %
            % ~phi_4 (x_2, x_3, x_5) for x_6 = 1
            \begin{tabular}{llll}
                \toprule
                $x_2$ & $x_3$ & $x_5$ & $\tilde{\phi}_4$\\
                \midrule
                0 & 0 & 0 & 132 \\
                1 & 0 & 0 & 144 \\
                0 & 1 & 0 & 216 \\
                1 & 1 & 0 & 108 \\
                0 & 0 & 1 & 132 \\
                1 & 0 & 1 & 168 \\
                0 & 1 & 1 & 120 \\
                1 & 1 & 1 & 60 \\
                \bottomrule
            \end{tabular}
            \hspace{3ex} \text{de modo que} \hspace{3ex}
            %
            % ~phi_{45} (x_2, x_3)
            \begin{tabular}{lllll}
                \toprule
                $x_2$ & $x_3$ & $\sum_{x_5} \tilde{\phi}_4(x_2, x_3, x_5)$ &  & $\tilde{\phi}_{45}$\\
                \midrule
                0 & 0 & 132 + 132 & = & 264 \\
                1 & 0 & 144 + 168 & = & 312 \\
                0 & 1 & 216 + 120 & = & 336 \\
                1 & 1 & 108 + 60  & = & 168 \\
                \bottomrule
            \end{tabular}
        \end{center}
        
        Finalmente encontre $\tilde{\phi}_{453}(x_2)$:
        \begin{center}
            %
            % phi_A (x_2, x_3) for x_1 = 0
            \begin{tabular}{lll}
                \toprule
                $x_2$ & $x_3$ & $\phi_A^{x_1=0}$\\
                \midrule
                0 & 0 & 4 \\
                1 & 0 & 2 \\
                0 & 1 & 2 \\
                1 & 1 & 6 \\
                \bottomrule
            \end{tabular}\hspace{2ex}
            %
            % ~phi_{45} (x_2, x_3)
            \begin{tabular}{lll}
                \toprule
                $x_2$ & $x_3$ & $\tilde{\phi}_{45}$\\
                \midrule
                0 & 0 & 264 \\
                1 & 0 & 312 \\
                0 & 1 & 336 \\
                1 & 1 & 168 \\
                \bottomrule
            \end{tabular}
        \end{center}
        de modo que:
        \begin{center}
            % ~phi_{453} (x_2)
            \begin{tabular}{llll}
                \toprule
                $x_2$ & $\sum_{x_3} \tilde{\phi}_{45}(x_2, x_3) \phi_A^{x_1=0}(x_2,x_3)$  &  & $\tilde{\phi}_{453}$\\
                \midrule
                0 & (4 * 264) + (2 * 336) & = & 1728 \\
                1 & (2 * 312) + (6 * 168) & = & 1632 \\
                \bottomrule
            \end{tabular}
        \end{center}\vspace{1ex}
        
        A constante de normaliza\'c\~ao \'e $Z = 1728 + 1632$. Nossa marginal condicional \'e portanto:
        \begin{equation}
            p(x_2 \mid x_1=0, x_6=1) = 
            \begin{pmatrix}
                1728 / Z \\
                1632 / Z \\
            \end{pmatrix} \approx 
            \begin{pmatrix}
                0.514 \\
                0.486 \\
            \end{pmatrix}
        \end{equation}
        
    \end{solution}
    
    %
    % Question c – find p(x_2 | x_1, x_6) with a good variable ordering.
    \item Agora determine $p(x_2 \mid x_1=0, x_6=1)$ com a ordem de elimina\'c\~ao $(x_5, x_4, x_3)$:
    
    \begin{exenumerate}
        \item Desenhe o grafo para $p(x_2, x_3, x_4, \mid x_1=0, x_6=1)$ ao marginalizar $x_5$ \\
        Calcule a tabela para o novo fator $\tilde{\phi}_{5}(x_4)$ \\
        
        \item Desenhe o grafo para $p(x_2, x_3 \mid x_1=0, x_6=1)$ ao marginalizar $x_4$ \\
        Calcule a tabela para o novo fator $\tilde{\phi}_{54}(x_2, x_3)$ \\
        
        \item Desenhe o grafo para $p(x_2 \mid x_1=0, x_6=1)$ ao marginalizar $x_3$ \\
        Calcule a tabela para o novo fator $\tilde{\phi}_{543}(x_2)$ 
    \end{exenumerate}
    
    %
    % Question c – solution
    \begin{solution}
        
        Come\'cando com o grafo de fatores para $p(x_2, x_3, x_4, x_5 \mid x_1=0, x_6=1)$:
        
        \begin{center}
            \begin{tikzpicture}[ugraph]
                \node[fact, label=above: $\phi_A^{x_1=0}$] (fa) at (1.5,0) {} ; %
                \node[cont] (x2) at (2.5,1)  {$x_2$} ; %
                \node[cont] (x3) at (2.5,-1)  {$x_3$} ; %
                \node[fact, label=above: $\phi_B$] (fb) at (3.5,0) {} ; %
                \node[cont] (x4) at (5,0)  {$x_4$} ; %
                \node[fact, label=above: $\phi_C$] (fc) at (6,1) {} ; %
                \node[cont] (x5) at (7.5,1)  {$x_5$} ; %
                \node[fact, label=right: $\phi_D^{x_6=1}$] (fd) at (6,-1) {} ; %
                
                \draw (fa) -- (x2);
                \draw (fa) -- (x3);
                \draw (x2) -- (fb);
                \draw (x3) -- (fb);
                \draw (fb) -- (x4);
                \draw (x4) -- (fc);
                \draw (x4) -- (fd);
                \draw (fc) -- (x5);
            \end{tikzpicture}
        \end{center}
        
        A marginaliza\'c\~ao de $x_5$ modifica o fator $\phi_C$:
        
        \begin{center}
            \begin{tikzpicture}[ugraph]
                \node[fact, label=above: $\phi_A^{x_1=0}$] (fa) at (1.5,0) {} ; %
                \node[cont] (x2) at (2.5,1)  {$x_2$} ; %
                \node[cont] (x3) at (2.5,-1)  {$x_3$} ; %
                \node[fact, label=above: $\phi_B$] (fb) at (3.5,0) {} ; %
                \node[cont] (x4) at (5,0)  {$x_4$} ; %
                \node[fact, label=right: $\tilde{\phi}_5$] (f5) at (6,1) {} ; %
                \node[fact, label=right: $\phi_D^{x_6=1}$] (fd) at (6,-1) {} ; %
                
                \draw (fa) -- (x2);
                \draw (fa) -- (x3);
                \draw (x2) -- (fb);
                \draw (x3) -- (fb);
                \draw (fb) -- (x4);
                \draw (x4) -- (f5);
                \draw (x4) -- (fd);
            \end{tikzpicture}
        \end{center}
        
        A marginaliza\'c\~ao de $x_4$ combina os tr\^es fatores $\phi_B$, $\tilde{\phi}_5$ e $\phi_D^{x_6=1}$:
        
        \begin{center}
            \begin{tikzpicture}[ugraph]
                \node[fact, label=above: $\phi_A^{x_1=0}$] (fa) at (1.5,0) {} ; %
                \node[cont] (x2) at (2.5,1)  {$x_2$} ; %
                \node[cont] (x3) at (2.5,-1)  {$x_3$} ; %
                \node[fact, label=right: $\tilde{\phi}_{54}$] (f54) at (3.5,0) {} ; %
                
                \draw (fa) -- (x2);
                \draw (fa) -- (x3);
                \draw (x2) -- (f54);
                \draw (x3) -- (f54);
            \end{tikzpicture}
        \end{center}
        
        A marginaliza\'c\~ao de $x_3$ combina os fatores $\phi_A^{x_1=0}$ e $\tilde{\phi}_{54}$:
        
        \begin{center}
            \begin{tikzpicture}[ugraph]
                \node[cont] (x2) at (0,0)  {$x_2$} ; %
                \node[fact, label=right: $\tilde{\phi}_{543}$] (f543) at (2,0) {} ; %
                
                \draw (x2) -- (f543);
            \end{tikzpicture}
        \end{center}
        
        Agora calculamos as tabelas para os novos fatores $\tilde{\phi}_5$, $\tilde{\phi}_{54}$, e $\tilde{\phi}_{543}$.
        
        Primeiro encontre $\tilde{\phi}_5(x_4)$:
        \begin{center}
            %
            % phi_C (x_4, x_5)
            \begin{tabular}{lll}
                \toprule
                $x_4$ & $x_5$ & $\phi_C$\\
                \midrule
                0 & 0 & 8 \\
                1 & 0 & 2 \\
                0 & 1 & 2 \\
                1 & 1 & 6 \\
                \bottomrule
            \end{tabular}
            \hspace{3ex} \text{de modo que} \hspace{3ex}
            %
            % ~phi_5 (x_4) for x_6 = 1
            \begin{tabular}{llll}
                \toprule
                $x_4$ & $\sum_{x_5} \phi_C(x_4, x_5)$ &  & $\tilde{\phi}_5$\\
                \midrule
                0 & 8 + 2 & = & 10 \\
                1 & 2 + 6 & = & 8 \\
                \bottomrule
            \end{tabular}
        \end{center}
        
        Em seguida, encontre $\tilde{\phi}_{54}(x_2, x_3)$:
        \begin{center}
            %
            % phi_B (x_2, x_3, x_4)
            \begin{tabular}{llll}
                \toprule
                $x_2$ & $x_3$ & $x_4$ & $\phi_B$\\
                \midrule
                0 & 0 & 0 & 2 \\
                1 & 0 & 0 & 2 \\
                0 & 1 & 0 & 4 \\
                1 & 1 & 0 & 2 \\
                0 & 0 & 1 & 6 \\
                1 & 0 & 1 & 8 \\
                0 & 1 & 1 & 4 \\
                1 & 1 & 1 & 2 \\
                \bottomrule
            \end{tabular}\hspace{2ex}
            %
            % ~phi_5 (x_4) for x_6 = 1
            \begin{tabular}{ll}
                \toprule
                $x_4$ & $\tilde{\phi}_5$\\
                \midrule
                0 & 10 \\
                1 & 8 \\
                \bottomrule
            \end{tabular}\hspace{2ex}
            %
            % phi_D (x_4) for x_6 = 1
            \begin{tabular}{ll}
                \toprule
                $x_4$ & $\phi_D^{x_6=1}$\\
                \midrule
                0 & 6 \\
                1 & 3 \\
                \bottomrule
            \end{tabular}
        \end{center}
        de modo que $\phi_*(x_2,x_3,x_4) = \phi_B(x_2, x_3, x_4) \tilde{\phi}_5(x_4) \phi_D^{x_6=1}(x_4)$ \'e igual a:
        %
        \begin{center}
            \begin{tabular}{lllll}
                \toprule
                $x_2$ & $x_3$ & $x_4$ & $\phi_*(x_2,x_3,x_4)$\\
                \midrule
                0 & 0 & 0 & 2 * 10 * 6 \\ 
                1 & 0 & 0 & 2 * 10 * 6 \\
                0 & 1 & 0 & 4 * 10 * 6 \\
                1 & 1 & 0 & 2 * 10 * 6 \\
                0 & 0 & 1 & 6 * 8 * 3 \\
                1 & 0 & 1 & 8 * 8 * 3 \\
                0 & 1 & 1 & 4 * 8 * 3 \\
                1 & 1 & 1 & 2 * 8 * 3 \\
                \bottomrule
            \end{tabular}
        \end{center}
        e:
        \begin{center}
            \begin{tabular}{lllll}
                \toprule
                $x_2$ & $x_3$ & $\sum_{x_4} \phi_B(x_2, x_3, x_4) \tilde{\phi}_5(x_4) \phi_D^{x_6=1}(x_4)$ &  & $\tilde{\phi}_{54}$\\
                \midrule
                0 & 0 & (2 * 10 * 6) + (6 * 8 * 3) & = & 264 \\
                1 & 0 & (2 * 10 * 6) + (8 * 8 * 3) & = & 312 \\
                0 & 1 & (4 * 10 * 6) + (4 * 8 * 3) & = & 336 \\
                1 & 1 & (2 * 10 * 6) + (2 * 8 * 3) & = & 168 \\
                \bottomrule
            \end{tabular}
        \end{center}
        
        Finalmente encontre $\tilde{\phi}_{543}(x_2)$:
        \begin{center}
            %
            % phi_A (x_2, x_3) for x_1 = 0
            \begin{tabular}{lll}
                \toprule
                $x_2$ & $x_3$ & $\phi_A^{x_1=0}$\\
                \midrule
                0 & 0 & 4 \\
                1 & 0 & 2 \\
                0 & 1 & 2 \\
                1 & 1 & 6 \\
                \bottomrule
            \end{tabular}\hspace{2ex}
            %
            % ~phi_{54} (x_2, x_3)
            \begin{tabular}{lll}
                \toprule
                $x_2$ & $x_3$ & $\tilde{\phi}_{54}$\\
                \midrule
                0 & 0 & 264 \\
                1 & 0 & 312 \\
                0 & 1 & 336 \\
                1 & 1 & 168 \\
                \bottomrule
            \end{tabular}
        \end{center}
        de modo que:
        \begin{center}
            %
            % ~phi_{ABCD}^{x_1=0, x_6=1} (x_2)
            \begin{tabular}{llll}
                \toprule
                $x_2$ & $\sum_{x_3} \tilde{\phi}_{54}(x_2, x_3) \phi_A^{x_1=0}(x_2,x_3)$ &  & $\tilde{\phi}_{543}$\\
                \midrule
                0 & (4 * 264) + (2 * 336) & = & 1728 \\
                1 & (2 * 312) + (6 * 168) & = & 1632 \\
                \bottomrule
            \end{tabular}
        \end{center}
        
        Assim como com a ordena\'c\~ao na parte anterior, devemos chegar ao
        mesmo resultado para nossa distribui\'c\~ao marginal condicional. A
        constante de normaliza\'c\~ao \'e $Z = 1728 + 1632$, portanto a
        marginal condicional \'e
        \begin{equation}
            p(x_2 \mid x_1=0, x_6=1) = 
            \begin{pmatrix}
                1728 / Z \\
                1632 / Z \\
            \end{pmatrix} \approx 
            \begin{pmatrix}
                0.514 \\
                0.486 \\
            \end{pmatrix}
        \end{equation}
        
    \end{solution}
    
    \item Qual ordem de vari\'aveis, $(x_4, x_5, x_3)$ ou $(x_5, x_4, x_3)$, voc\^e prefere?
    
    \begin{solution}
        A ordena\'c\~ao $(x_5, x_4, x_3)$ \'e mais barata computacionalmente e deve ser preferida
        em rela\'c\~ao \`a ordena\'c\~ao $(x_4, x_5, x_3)$.
        
        A raz\~ao para a diferen\'ca no custo \'e que $x_4$ possui tr\^es
        vizinhos no grafo de fatores para $p(x_2, x_3, x_4, x_5 \mid
        x_1=0, x_6=1)$. No entanto, ap\'os a elimina\'c\~ao de $x_5$, que possui
        apenas um vizinho, $x_4$ fica com apenas dois vizinhos restantes. 
        Eliminar vari\'aveis com mais vizinhos leva a
        fatores (tempor\'arios) maiores e, consequentemente, a um custo maior. Podemos ver
        isso nas tabelas que foram geradas durante o c\'alculo (ou nos n\'umeros
        que precisaram ser somados): para a ordena\'c\~ao $(x_4, x_5, x_3)$, a maior tabela teve $2^4$ entradas, enquanto para $(x_5, x_4, x_3)$, ela teve $2^3$ entradas.
        
        Escolher uma ordena\'c\~ao razo\'avel de vari\'aveis tem um efeito direto na
        complexidade computacional da elimina\'c\~ao de vari\'aveis. Esse efeito
        torna-se ainda mais pronunciado quando o dom\'inio de nossas vari\'aveis discretas
        tem um tamanho maior que 2 (vari\'aveis bin\'arias), ou se as
        vari\'aveis forem cont\'inuas.
        
        \begin{center}
            \begin{tikzpicture}[ugraph]
                \node[fact, label=above: $\phi_A^{x_1=0}$] (fa) at (1.5,0) {} ; %
                \node[cont] (x2) at (2.5,1)  {$x_2$} ; %
                \node[cont] (x3) at (2.5,-1)  {$x_3$} ; %
                \node[fact, label=above: $\phi_B$] (fb) at (3.5,0) {} ; %
                \node[cont] (x4) at (5,0)  {$x_4$} ; %
                \node[fact, label=above: $\phi_C$] (fc) at (6,1) {} ; %
                \node[cont] (x5) at (7.5,1)  {$x_5$} ; %
                \node[fact, label=right: $\phi_D^{x_6=1}$] (fd) at (6,-1) {} ; %
                
                \draw (fa) -- (x2);
                \draw (fa) -- (x3);
                \draw (x2) -- (fb);
                \draw (x3) -- (fb);
                \draw (fb) -- (x4);
                \draw (x4) -- (fc);
                \draw (x4) -- (fd);
                \draw (fc) -- (x5);
            \end{tikzpicture}
        \end{center}
        
    \end{solution}
    
\end{exenumerate}



\ex{Escolha da ordem de elimina\'c\~ao em grafos de fatores}

Gostar\'iamos de calcular a marginal $p(x_1)$ por elimina\'c\~ao de vari\'aveis
para uma fmp conjunta representada pelo seguinte grafo de
fatores. Todas as vari\'aveis $x_i$ podem assumir $K$ valores diferentes. \\
\begin{center}
    \scalebox{1}{ % x,y
        \begin{tikzpicture}[ugraph]
            \node[cont] (x1) at (0,0) {$x_1$};
            \node[fact, label=above: $\phi_a$] (fa) at (1,-1) {};
            \node[cont] (x2) at (2,0) {$x_2$};
            \node[fact, label={[xshift=0.15cm,  yshift=0cm]left: $\phi_b$}] (fb) at (2,-1) {};
            \node[cont] (x3) at (4,0) {$x_3$};
            \node[fact, label=above: $\phi_c$] (fc) at (3,-1) {};
            
            \node[cont] (x4) at (2,-2) {$x_4$};
            
            \node[cont] (x5) at (0,-4) {$x_5$};
            \node[fact, label=above: $\phi_d$] (fd) at (1,-3) {};
            \node[cont] (x6) at (2,-4) {$x_6$};
            \node[fact, label={[xshift=0.15cm,  yshift=0cm]left: $\phi_e$}] (fe) at (2,-3) {};
            \node[cont] (x7) at (4,-4) {$x_7$};
            \node[fact, label=above: $\phi_f$] (ff) at (3,-3) {};
            
            \draw(x1) -- (fa);
            \draw(x2) -- (fb);
            \draw(x3) -- (fc);
            \draw(x4) -- (fa);
            \draw(x4) -- (fb);
            \draw(x4) -- (fc);
            
            \draw(x4) -- (fc);
            
            \draw(x4) -- (fd);
            \draw(x4) -- (fe);
            \draw(x4) -- (ff);
            
            \draw(fd) -- (x5);
            \draw(fe) -- (x6);
            \draw(ff) -- (x7);
    \end{tikzpicture}}
\end{center}
\begin{exenumerate}
    \item Um amigo prop\~oe a ordem de elimina\'c\~ao $x_4, x_5, x_6, x_7, x_3, x_2$, i.e.\ fazer $x_4$ primeiro e $x_2$ por \'ultimo. Explique por que isso \'e computacionalmente ineficiente.
    
    \begin{solution}
        
        De acordo com o grafo de fatores, $p(x_1, \ldots, x_7)$ se fatoriza como
        \begin{align}
            p(x_1, \ldots, x_7) &\propto \phi_a(x_1,x_4)\phi_b(x_2,x_4)\phi_c(x_3,x_4)\phi_d(x_5,x_4)\phi_e(x_6,x_4)\phi_f(x_7,x_4)
        \end{align}
        Se optarmos por eliminar $x_4$ primeiro, i.e. calcular
        \begin{align}
            p(x_1,x_2,x_3,x_5, x_6, x_7) &= \sum_{x_4}   p(x_1, \ldots, x_7) \\
            & \propto \sum_{x_4} \phi_a(x_1,x_4)\phi_b(x_2,x_4)\phi_c(x_3,x_4)\phi_d(x_5,x_4)\phi_e(x_6,x_4)\phi_f(x_7,x_4)
        \end{align}
        n\~ao podemos retirar nenhum dos fatores da soma, j\'a que cada um deles depende de $x_4$. Isso significa que o custo para somar $x_4$ para todas as combina\'c\~oes das seis vari\'aveis $(x_1,x_2,x_3,x_5, x_6, x_7)$ \'e $K^7$. Al\'em disso, o novo fator
        \begin{equation}
            \tilde{\phi}(x_1,x_2,x_3,x_5, x_6, x_7) =  \sum_{x_4} \phi_a(x_1,x_4)\phi_b(x_2,x_4)\phi_c(x_3,x_4)\phi_d(x_5,x_4)\phi_e(x_6,x_4)\phi_f(x_7,x_4)
        \end{equation}
        n\~ao se fatoriza mais, de modo que as elimina\'c\~oes de vari\'aveis subsequentes tamb\'em ser\~ao caras.
    \end{solution}
    
    \item Proponha uma ordem de elimina\'c\~ao que alcance o custo computacional $O(K^2)$ por elimina\'c\~ao de vari\'avel e explique por que ela o faz.
    
    \begin{solution}
        
        Qualquer ordena\'c\~ao onde $x_4$ seja eliminado por \'ultimo servir\'a. Em qualquer est\'agio, a elimina\'c\~ao de uma das vari\'aveis $x_2, x_3, x_5, x_6, x_7$ \'e ent\~ao uma opera\'c\~ao $O(K^2)$. Isso ocorre porque, e.g.:
        \begin{align}
            p(x_1, \ldots, x_6) & = \sum_{x_7} p(x_1, \ldots, x_7)\\
            & \propto \phi_a(x_1,x_4)\phi_b(x_2,x_4)\phi_c(x_3,x_4)\phi_d(x_5,x_4)\phi_e(x_6,x_4) \underbrace{\sum_{x_7} \phi_f(x_7,x_4)}_{\tilde{\phi}_7(x_4)}\\
            & \propto \phi_a(x_1,x_4)\phi_b(x_2,x_4)\phi_c(x_3,x_4)\phi_d(x_5,x_4)\phi_e(x_6,x_4) \tilde{\phi}_7(x_4)
        \end{align}
        onde o c\'alculo de $\tilde{\phi}_7(x_4)$ para todos os valores de $x_4$ \'e $O(K^2)$. Al\'em disso,
        \begin{align}
            p(x_1, \ldots, x_5) & = \sum_{x_6} p(x_1, \ldots, x_6)\\
            & \propto \phi_a(x_1,x_4)\phi_b(x_2,x_4)\phi_c(x_3,x_4)\phi_d(x_5,x_4) \tilde{\phi}_7(x_4) \sum_{x_6}\phi_e(x_6,x_4)\\
            & \propto \phi_a(x_1,x_4)\phi_b(x_2,x_4)\phi_c(x_3,x_4)\phi_d(x_5,x_4) \tilde{\phi}_7(x_4) \tilde{\phi}_6(x_4), 
        \end{align}
        onde o c\'alculo de $\tilde{\phi}_6(x_4)$ para todos os valores de $x_4$ \'e novamente $O(K^2)$. Continuando dessa maneira, obt\'em-se:
        \begin{align}
            p(x_1, x_4) &\propto \phi_a(x_1,x_4) \tilde{\phi}_2(x_4)\tilde{\phi}_3(x_4)\tilde{\phi}_5(x_4)\tilde{\phi}_6(x_4)\tilde{\phi}_7(x_4).
        \end{align}
        onde cada fator derivado $\tilde{\phi}$ tem custo $O(K^2)$. A elimina\'c\~ao de $x_4$ e a normaliza\'c\~ao da fmp tamb\'em \'e uma opera\'c\~ao $O(K^2)$.
    \end{solution}
    
\end{exenumerate}